% 参考文献：由底本（Springer LNM 2386）文字层转录，条目保持英文原貌；
% 编号即原文引用编号。DOI/arXiv 间距已做机器清理，个别符号可能与印刷本略有差异。
\chapter*{参考文献}
\addcontentsline{toc}{chapter}{参考文献}
\begingroup
\small
\setlength{\parindent}{0pt}
\setlength{\parskip}{3pt}
\noindent\hangindent=2.2em\hangafter=1 [1]\ S. L. Adler. “Axial vector vertex in spinor electrodynamics.” Phys. Rev. 177 (1969), pp. 2426--2438. doi: 10.1103/PhysRev.177.2426.\par
\noindent\hangindent=2.2em\hangafter=1 [2]\ S. L. Adler. “Anomalies to all orders.” 50 years of Yang-Mills theory. Ed. by G. ’t Hooft. 2005, pp. 187--228. doi: 10.1142/9789812567147\_0009.arXiv: hep-th/0405040.\par
\noindent\hangindent=2.2em\hangafter=1 [3]\ I. Affleck, M. Dine, and N. Seiberg. “Supersymmetry Breaking by Instantons.” Phys. Rev. Lett. 51 (1983), p. 1026. doi: 10.1103/PhysRevLett.51.1026.\par
\noindent\hangindent=2.2em\hangafter=1 [4]\ I. Affleck, M. Dine, and N. Seiberg. “Dynamical Supersymmetry Breaking in Supersymmetric QCD.” Nucl. Phys. B 241 (1984), pp. 493--534. doi: 10.1016/0550-3213(84)90058-0.\par
\noindent\hangindent=2.2em\hangafter=1 [5]\ I. Affleck, M. Dine, and N. Seiberg. “Dynamical Supersymmetry Breaking in Four- Dimensions and Its Phenomenological Implications.” Nucl. Phys. B 256 (1985), pp. 557-- 599. doi: 10.1016/0550-3213(85)90408-0.\par
\noindent\hangindent=2.2em\hangafter=1 [6]\ O. Aharony, S. S. Gubser, J. M. Maldacena, H. Ooguri, and Y. Oz. “Large N field theories, string theory and gravity.” Phys. Rept. 323 (2000), pp. 183--386. doi: 10.1016/S0370-1573(99)00083-6.arXiv: hep-th/9905111.\par
\noindent\hangindent=2.2em\hangafter=1 [7]\ O. Aharony, N. Seiberg, and Y. Tachikawa. “Reading between the lines of four-dimensional gauge theories.” JHEP 08 (2013), p. 115. doi: 10.1007/JHEP08(2013)115.arXiv: 1305.0318 [hep-th].\par
\noindent\hangindent=2.2em\hangafter=1 [8]\ O. Aharony, Y. Tachikawa, and K. Gomi. “S-folds and 4d N=3 superconformal field theories.” JHEP 06 (2016), p. 044. doi: 10.1007/JHEP06(2016)044.arXiv: 1602.08638 [hep-th].\par
\noindent\hangindent=2.2em\hangafter=1 [9]\ M. Akhond, G. Arias-Tamargo, A. Mininno, H.-Y. Sun, Z. Sun, Y. Wang, and F. Xu. “The hitchhiker’s guide to 4d N = 2 superconformal field theories.” SciPost Phys. Lect. Notes 64 (2022), p. 1. doi: 10.21468/SciPostPhysLectNotes.64.arXiv: 2112.14764 [hep-th].\par
\noindent\hangindent=2.2em\hangafter=1 [10]\ L. F. Alday, D. Gaiotto, S. Gukov, Y. Tachikawa, and H. Verlinde. “Loop and surface operators in N=2 gauge theory and Liouville modular geometry.” JHEP 01 (2010), p. 113. doi: 10.1007/JHEP01(2010)113.arXiv: 0909.0945 [hep-th].\par
\noindent\hangindent=2.2em\hangafter=1 [11]\ L. F. Alday, D. Gaiotto, and Y. Tachikawa. “Liouville Correlation Functions from Fourdimensional Gauge Theories.” Lett. Math. Phys. 91 (2010), pp. 167--197. doi: 10.1007/s11005-010-0369-5.arXiv: 0906.3219 [hep-th].\par
\noindent\hangindent=2.2em\hangafter=1 [12]\ D. Alessandrini. Introduction to Spectral Networks. Mini-course at University of Illinois at Urbana-Champaign. 2014.\par
\noindent\hangindent=2.2em\hangafter=1 [13]\ D. Alessandrini, A. Berenstein, V. Retakh, E. Rogozinnikov, and A. Wienhard. “Symplectic groups over noncommutative algebras.” Selecta Mathematica 28.4 (2022), p. 82. doi: 10.1007/s00029-022-00787-x.arXiv: 2106.08736 [math.DG].\par
\noindent\hangindent=2.2em\hangafter=1 [14]\ D. Alessandrini, O. Guichard, E. Rogozinnikov, and A. Wienhard. Noncommutative coordinates for symplectic representations. Vol. 1504. Mem. Am. Math. Soc. Providence, RI: American Mathematical Society (AMS), 2024. isbn: 978-1-4704-7930-5. doi: 10.1090/memo/1504.arXiv: 1911.08014 [math.DG].\par
\noindent\hangindent=2.2em\hangafter=1 [15]\ M. Alim, S. Cecotti, C. Cordova, S. Espahbodi, A. Rastogi, and C. Vafa. “BPS Quivers and Spectra of Complete N=2 Quantum Field Theories.” Commun. Math. Phys. 323 (2013), pp. 1185--1227. doi: 10.1007/s00220-013-1789-8.arXiv: 1109.4941 [hep-th].\par
\noindent\hangindent=2.2em\hangafter=1 [16]\ M. Alim, S. Cecotti, C. Cordova, S. Espahbodi, A. Rastogi, and C. Vafa. “N = 2 quantum field theories and their BPS quivers.” Adv. Theor. Math. Phys. 18.1 (2014), pp. 27--127. doi: 10.4310/ATMP.2014.v18.n1.a2.arXiv: 1112.3984 [hep-th].\par
\noindent\hangindent=2.2em\hangafter=1 [17]\ M. Alim, L. Hollands, and I. Tulli. “Quantum Curves, Resurgence and Exact WKB.” SIGMA 19 (2023), p. 009. doi: 10.3842/SIGMA.2023.009.arXiv: 2203.08249 [hep-th].\par
\noindent\hangindent=2.2em\hangafter=1 [18]\ L. Alvarez-Gaume and S. F. Hassan. “Introduction to S duality in N=2 supersymmetric gauge theories: A Pedagogical review of the work of Seiberg and Witten.” Fortsch. Phys. 45 (1997), pp. 159--236. doi: 10.1002/prop.2190450302.arXiv: hep-th/9701069.\par
\noindent\hangindent=2.2em\hangafter=1 [19]\ L. Alvarez-Gaume and M. A. Vazquez-Mozo. “Introductory lectures on quantum field theory.” 3rd CERN-CLAF School of High Energy Physics. 2006, pp. 1--80. doi: 10.5170/CERN-2006-015.1.arXiv: hep-th/0510040.\par
\noindent\hangindent=2.2em\hangafter=1 [20]\ T. Ando. “Totally positive matrices.” Linear algebra and its applications 90 (1987), pp. 165--219.\par
\noindent\hangindent=2.2em\hangafter=1 [21]\ R. Argurio, G. Ferretti, and R. Heise. “An Introduction to supersymmetric gauge theories and matrix models.” Int. J. Mod. Phys. A 19 (2004), pp. 2015--2078. doi: 10.1142/S0217751X04018038.arXiv: hep-th/0311066.\par
\noindent\hangindent=2.2em\hangafter=1 [22]\ P. C. Argyres. An Introduction to Global Supersymmetry. University of Cornell Lecture Notes. 2001.\par
\noindent\hangindent=2.2em\hangafter=1 [23]\ P. C. Argyres and A. Buchel. “New S dualities in N=2 supersymmetric SU(2) x SU(2) gauge theory.” JHEP 11 (1999), p. 014. doi: 10.1088/1126-6708/1999/11/014.arXiv: hep-th/9910125.\par
\noindent\hangindent=2.2em\hangafter=1 [24]\ P. C. Argyres and M. R. Douglas. “New phenomena in SU(3) supersymmetric gauge theory.” Nucl. Phys. B 448 (1995), pp. 93--126. doi: 10.1016/0550-3213(95)00281-V.arXiv: hep-th/9505062.\par
\noindent\hangindent=2.2em\hangafter=1 [25]\ P. C. Argyres and A. E. Faraggi. “The vacuum structure and spectrum of N=2 supersymmetric SU(n) gauge theory.” Phys. Rev. Lett. 74 (1995), pp. 3931--3934. doi: 10.1103/PhysRevLett.74.3931.arXiv: hep-th/9411057.\par
\noindent\hangindent=2.2em\hangafter=1 [26]\ P. C. Argyres, M. R. Plesser, and N. Seiberg. “The Moduli space of vacua of N=2 SUSY QCD and duality in N=1 SUSY QCD.” Nucl. Phys. B 471 (1996), pp. 159--194. doi: 10.1016/0550-3213(96)00210-6.arXiv: hep-th/9603042.\par
\noindent\hangindent=2.2em\hangafter=1 [27]\ P. C. Argyres, M. R. Plesser, N. Seiberg, and E. Witten. “New N=2 superconformal field theories in four-dimensions.” Nucl. Phys. B 461 (1996), pp. 71--84. doi: 10.1016/0550-3213(95)00671-0.arXiv: hep-th/9511154.\par
\noindent\hangindent=2.2em\hangafter=1 [28]\ P. C. Argyres, M. R. Plesser, and A. D. Shapere. “The Coulomb phase of N=2 supersymmetric QCD.” Phys. Rev. Lett. 75 (1995), pp. 1699--1702. doi: 10.1103/PhysRevLett.75.1699.arXiv: hep-th/9505100.\par
\noindent\hangindent=2.2em\hangafter=1 [29]\ P. C. Argyres and N. Seiberg. “S-duality in N=2 supersymmetric gauge theories.” JHEP 12 (2007), p. 088. doi: 10.1088/1126-6708/2007/12/088.arXiv: 0711.0054 [hep-th].\par
\noindent\hangindent=2.2em\hangafter=1 [30]\ P. C. Argyres and A. D. Shapere. “The Vacuum structure of N=2 superQCD with classical gauge groups.” Nucl. Phys. B 461 (1996), pp. 437--459. doi: 10.1016/0550-3213(95)00661-3.arXiv: hep-th/9509175.\par
\noindent\hangindent=2.2em\hangafter=1 [31]\ N. Arkani-Hamed and H. Murayama. “Holomorphy, rescaling anomalies and exact beta functions in supersymmetric gauge theories.” JHEP 06 (2000), p. 030. doi: 10.1088/1126-6708/2000/06/030.arXiv: hep-th/9707133.\par
\noindent\hangindent=2.2em\hangafter=1 [32]\ V. I. Arnold. Mathematical methods of classical mechanics. Translated by K. Vogtman and A. Weinstein. Vol. 60. Grad. Texts Math. Springer, Cham, 1978.\par
\noindent\hangindent=2.2em\hangafter=1 [33]\ M. F. Atiyah and R. Bott. “The Yang-Mills equations over Riemann surfaces.” Phil. Trans. Roy. Soc. Lond. A 308 (1982), pp. 523--615.\par
\noindent\hangindent=2.2em\hangafter=1 [34]\ M. F. Atiyah and R. S. Ward. “Instantons and Algebraic Geometry.” Commun. Math. Phys. 55 (1977), pp. 117--124. doi: 10.1007/BF01626514.\par
\noindent\hangindent=2.2em\hangafter=1 [35]\ M. Atiyah. “Topological quantum field theories.” Publ. Math., Inst. Hautes Étud. Sci. 68 (1988), pp. 175--186. issn: 0073-8301. doi: 10.1007/BF02698547.\par
\noindent\hangindent=2.2em\hangafter=1 [36]\ M. F. Atiyah, N. J. Hitchin, and I. M. Singer. “Self-duality in four-dimensional Riemannian geometry.” Proceedings of the Royal Society of London. A. Mathematical and Physical Sciences 362.1711 (1978), pp. 425--461.\par
\noindent\hangindent=2.2em\hangafter=1 [37]\ J. Baez and J. P. Muniain. Gauge fields, knots and gravity. Series on Knots and Everything Vol. 4. World Scientific, 1995.\par
\noindent\hangindent=2.2em\hangafter=1 [38]\ S. Banerjee, P. Longhi, and M. Romo. “Exploring 5d BPS Spectra with Exponential Networks.” Annales Henri Poincare 20.12 (2019), pp. 4055--4162. doi: 10.1007/s00023-019-00851-x.arXiv: 1811.02875 [hep-th].\par
\noindent\hangindent=2.2em\hangafter=1 [39]\ S. Banerjee, P. Longhi, and M. Romo. “Exponential BPS Graphs and D Brane Counting on Toric Calabi-Yau Threefolds: Part I.” Commun. Math. Phys. 388.2 (2021), pp. 893--945. doi: 10.1007/s00220-021-04242-4.arXiv: 1910.05296 [hep-th].\par
\noindent\hangindent=2.2em\hangafter=1 [40]\ S. Banerjee, P. Longhi, and M. Romo. “A-branes, Foliations and Localization.” Annales Henri Poincare 24.4 (2023), pp. 1077--1136. doi: 10.1007/s00023-022-01231-8.arXiv: 2201.12223 [hep-th].\par
\noindent\hangindent=2.2em\hangafter=1 [41]\ S. Banerjee, P. Longhi, and M. Romo. “Exponential BPS Graphs and D-Brane Counting on Toric Calabi--Yau Threefolds: Part II.” Commun. Math. Phys. 406.7 (2025), p. 172. doi: 10.1007/s00220-025-05333-2.arXiv: 2012.09769 [hep-th].\par
\noindent\hangindent=2.2em\hangafter=1 [42]\ S. Banerjee, M. Romo, R. Senghaas, and J. Walcher. “Exponential Networks for Linear Partitions” (Mar. 2024). arXiv: 2403.14588 [hep-th].\par
\noindent\hangindent=2.2em\hangafter=1 [43]\ K. Becker, M. Becker, and J. H. Schwarz. String theory and M-theory: A modern introduction. Cambridge University Press, Dec. 2006. isbn: 978-0-511-81608-6. doi: 10.1017/CBO9780511816086.\par
\noindent\hangindent=2.2em\hangafter=1 [44]\ A. Beilinson and V. Drinfeld. “Opers” (2005). arXiv: math/0501398.\par
\noindent\hangindent=2.2em\hangafter=1 [45]\ A. A. Belavin, A. M. Polyakov, A. S. Schwartz, and Y. S. Tyupkin. “Pseudoparticle Solutions of the Yang-Mills Equations.” Phys. Lett. B 59 (1975). Ed. by J. C. Taylor, pp. 85--87. doi: 10.1016/0370-2693(75)90163-X.\par
\noindent\hangindent=2.2em\hangafter=1 [46]\ J. S. Bell and R. Jackiw. “A PCAC puzzle: $\pi$0 →$\gamma$$\gamma$in the $\sigma$model.” Nuovo Cim. A 60 (1969), pp. 47--61. doi: 10.1007/BF02823296.\par
\noindent\hangindent=2.2em\hangafter=1 [47]\ R. Benedetti and C. Petronio. Lectures on hyperbolic geometry. Universitext. Berlin etc.: Springer-Verlag, 1992. isbn: 3-540-55534-X.\par
\noindent\hangindent=2.2em\hangafter=1 [48]\ A. Berenstein, S. Fomin, and A. Zelevinsky. “Cluster algebras. III: Upper bounds and double Bruhat cells.” Duke Math. J. 126.1 (2005), pp. 1--52. issn: 0012-7094. doi: 10.1215/S0012-7094-04-12611-9.arXiv: math/0305434.\par
\noindent\hangindent=2.2em\hangafter=1 [49]\ D. Berenstein and M. R. Douglas. “Seiberg duality for quiver gauge theories” (July 2002). arXiv: hep-th/0207027.\par
\noindent\hangindent=2.2em\hangafter=1 [50]\ M. Bertolini. Supersymmetry. From the Basics to Exact Results in Gauge Theories. Lecture Notes in Physics Vol. 86. World Scientific, 2024. isbn: 978-981-98-0024-7. doi: 10.1142/14026.\par
\noindent\hangindent=2.2em\hangafter=1 [51]\ J. Beyrer, O. Guichard, F. Labourie, B. Pozzetti, and A. Wienhard. “Positivity, cross-ratios and the Collar Lemma” (2024). arXiv: 2409.06294 [math.DG].\par
\noindent\hangindent=2.2em\hangafter=1 [52]\ J. Beyrer and M. B. Pozzetti. “Positive surface group representations in PO( p, q).” J. Eur. Math. Soc. (JEMS) 28.6 (2026), pp. 2307--2349. issn: 1435-9855. doi: 10.4171/jems/1560.\par
\noindent\hangindent=2.2em\hangafter=1 [53]\ L. Bhardwaj, L. E. Bottini, L. Fraser-Taliente, L. Gladden, D. S. W. Gould, A. Platschorre, and H. Tillim. “Lectures on generalized symmetries.” Phys. Rept. 1051 (2024), pp. 1--87. doi: 10.1016/j.physrep.2023.11.002.arXiv: 2307.07547 [hep-th].\par
\noindent\hangindent=2.2em\hangafter=1 [54]\ M. Bianchi, S. Kovacs, and G. Rossi. “Instantons and Supersymmetry.” Lect. Notes Phys. 737 (2008), pp. 303--470. arXiv: hep-th/0703142.\par
\noindent\hangindent=2.2em\hangafter=1 [55]\ A. Bilal, V. V. Fock, and I. I. Kogan. “On the origin of W algebras.” Nucl. Phys. B 359 (1991), pp. 635--672. doi: 10.1016/0550-3213(91)90075-9.\par
\noindent\hangindent=2.2em\hangafter=1 [56]\ A. Bilal. “Introduction to supersymmetry” (Jan. 2001). arXiv: hep-th/0101055.\par
\noindent\hangindent=2.2em\hangafter=1 [57]\ A. Bilal. “Lectures on Anomalies” (Feb. 2008). arXiv: 0802.0634 [hep-th].\par
\noindent\hangindent=2.2em\hangafter=1 [58]\ A. Bilal and F. Ferrari. “Curves of marginal stability, and weak and strong coupling BPS spectra in N=2 supersymmetric QCD.” Nucl. Phys. B 480 (1996), pp. 589--622. doi: 10.1016/S0550-3213(96)00480-4.arXiv: hep-th/9605101.\par
\noindent\hangindent=2.2em\hangafter=1 [59]\ E. B. Bogomolny. “Stability of Classical Solutions.” Sov. J. Nucl. Phys. 24 (1976), p. 449.\par
\noindent\hangindent=2.2em\hangafter=1 [60]\ F. Bonahon and H. Wong. “Quantum traces for representations of surface groups in SL(2,C).” Geom. Topol. 15.3 (2011), pp. 1569--1615. issn: 1465-3060. doi: 10.2140/gt.2011.15.1569.\par
\noindent\hangindent=2.2em\hangafter=1 [61]\ N. Bourbaki. Eléments de mathématique. Groupes et algèbres de Lie. Chapitres 7 et 8. Reprint of the 1975 original. Berlin: Springer, 2006. isbn: 3-540-33939-6. doi: 10.1007/978-3-540-33977-9.\par
\noindent\hangindent=2.2em\hangafter=1 [62]\ S. Bradlow. “Global properties of Higgs bundle moduli spaces.” Moduli spaces and vector bundles -- new trends. VBAC 2022 conference in honor of Peter Newstead’s 80th birthday, University of Warwick, Coventry, United kingdom, July 25--29, 2022. Providence, RI: American Mathematical Society (AMS), 2024, pp. 39--77. isbn: 978-1-4704-7646-5. doi: 10.1090/conm/803/16094.arXiv: 2312.00762 [math.AG].\par
\noindent\hangindent=2.2em\hangafter=1 [63]\ L. Bramley, L. Hollands, and S. Murugesan. Les Houches lectures on non-perturbative Seiberg-Witten geometry. Preprint, arXiv: 2503.21742 [hep-th] (2025). 2025.\par
\noindent\hangindent=2.2em\hangafter=1 [64]\ A. Brandhuber and K. Landsteiner. “On the monodromies of N=2 supersymmetric Yang- Mills theory with gauge group SO(2n).” Phys. Lett. B 358 (1995), pp. 73--80. doi: 10.1016/0370-2693(95)00986-U.arXiv: hep-th/9507008.\par
\noindent\hangindent=2.2em\hangafter=1 [65]\ A. Brandhuber and S. Stieberger. “Selfdual strings and stability of BPS states in N=2 SU(2) gauge theories.” Nucl. Phys. B 488 (1997), pp. 199--222. doi: 10.1016/S0550-3213(96)00656-6.arXiv: hep-th/9610053.\par
\noindent\hangindent=2.2em\hangafter=1 [66]\ T. Bridgeland and I. Smith. “Quadratic differentials as stability conditions.” Publ. Math., Inst. Hautes Étud. Sci. 121 (2015), pp. 155--278. issn: 0073-8301. doi: 10.1007/s10240-014-0066-5.arXiv: 1302.7030 [math.AG].\par
\noindent\hangindent=2.2em\hangafter=1 [67]\ M. Bridgeman, R. Canary, and A. Sambarino. “An introduction to pressure metrics for higher Teichmüller spaces.” Ergodic Theory Dyn. Syst. 38.6 (2018), pp. 2001--2035. issn: 0143-3857. doi: 10.1017/etds.2016.111.arXiv: 1507.03894 [math.DG].\par
\noindent\hangindent=2.2em\hangafter=1 [68]\ L. Brillouin. “La mécanique ondulatoire de Schrödinger; une méthode générale de résolution par approximations successives.” Compt. Rend. Hebd. Seances Acad. Sci. 183 (1926), pp. 24--26.\par
\noindent\hangindent=2.2em\hangafter=1 [69]\ M. Burger, A. Iozzi, F. Labourie, and A. Wienhard. “Maximal representations of surface groups: symplectic Anosov structures.” Pure Appl. Math. Q. 1.3 (2005), pp. 543--590. issn: 1558-8599. doi: 10.4310/PAMQ.2005.v1.n3.a5.arXiv: math/0506079.\par
\noindent\hangindent=2.2em\hangafter=1 [70]\ M. Burger, A. Iozzi, and A. Wienhard. “Surface group representations with maximal Toledo invariant.” C. R., Math., Acad. Sci. Paris 336.5 (2003), pp. 387--390. issn: 1631- 073X. doi: 10.1016/S1631-073X(03)00065-7.arXiv: math/0605656.\par
\noindent\hangindent=2.2em\hangafter=1 [71]\ Y. M. Burman. “Triangulations of surfaces with boundary and the homotopy principle for functions without critical points.” Ann. Global Anal. Geom. 17.3 (1999), pp. 221--238. issn: 0232-704X. doi: 10.1023/A:1006556632099.\par
\noindent\hangindent=2.2em\hangafter=1 [72]\ P. Candelas, G. T. Horowitz, A. Strominger, and E. Witten. “Vacuum configurations for superstrings.” Nucl. Phys. B 258 (1985), pp. 46--74. doi: 10.1016/0550-3213(85)90602-9.\par
\noindent\hangindent=2.2em\hangafter=1 [73]\ S. Cecotti. “The quiver approach to the BPS spectrum of a 4d N=2 gauge theory.” Proc. Symp. Pure Math. 90 (2015). Ed. by R. Donagi, S. Katz, A. Klemm, and D. R. Morrison, pp. 3--18. arXiv: 1212.3431 [hep-th].\par
\noindent\hangindent=2.2em\hangafter=1 [74]\ S. Cecotti, P. Fendley, K. A. Intriligator, and C. Vafa. “A New supersymmetric index.” Nucl. Phys. B 386 (1992), pp. 405--452. doi: 10.1016/0550-3213(92)90572-S.arXiv: hep-th/9204102.\par
\noindent\hangindent=2.2em\hangafter=1 [75]\ S. Cecotti and C. Vafa. “On classification of N=2 supersymmetric theories.” Commun. Math. Phys. 158 (1993), pp. 569--644. doi: 10.1007/BF02096804.arXiv: hep-th/9211097.\par
\noindent\hangindent=2.2em\hangafter=1 [76]\ O. Chacaltana and J. Distler. “Tinkertoys for Gaiotto Duality.” JHEP 11 (2010), p. 099. doi: 10.1007/JHEP11(2010)099.arXiv: 1008.5203 [hep-th].\par
\noindent\hangindent=2.2em\hangafter=1 [77]\ O. Chacaltana and J. Distler. “Tinkertoys for the DNseries.” JHEP 02 (2013), p. 110. doi: 10.1007/JHEP02(2013)110.arXiv: 1106.5410 [hep-th].\par
\noindent\hangindent=2.2em\hangafter=1 [78]\ O. Chacaltana, J. Distler, and Y. Tachikawa. “Nilpotent orbits and codimension-two defects of 6d N=(2,0) theories.” Int. J. Mod. Phys. A 28 (2013), p. 1340006. doi: 10.1142/S0217751X1340006X.arXiv: 1203.2930 [hep-th].\par
\noindent\hangindent=2.2em\hangafter=1 [79]\ O. Chacaltana, J. Distler, and Y. Tachikawa. “Gaiotto duality for the twisted A2N−1 series.” JHEP 05 (2015), p. 075. doi: 10.1007/JHEP05(2015)075.arXiv: 1212.3952 [hep-th].\par
\noindent\hangindent=2.2em\hangafter=1 [80]\ O. Chacaltana, J. Distler, and A. Trimm. “Tinkertoys for the E6 theory.” JHEP 09 (2015), p. 007. doi: 10.1007/JHEP09(2015)007.arXiv: 1403.4604 [hep-th].\par
\noindent\hangindent=2.2em\hangafter=1 [81]\ O. Chacaltana, J. Distler, and A. Trimm. “Tinkertoys for the Twisted D-Series.” JHEP 04 (2015), p. 173. doi: 10.1007/JHEP04(2015)173.arXiv: 1309.2299 [hep-th].\par
\noindent\hangindent=2.2em\hangafter=1 [82]\ O. Chacaltana, J. Distler, and A. Trimm. “Tinkertoys for the Twisted E6 Theory” (Jan. 2015). arXiv: 1501.00357 [hep-th].\par
\noindent\hangindent=2.2em\hangafter=1 [83]\ O. Chacaltana, J. Distler, and A. Trimm. “Tinkertoys for the Z3-twisted D4 Theory” (Jan. 2016). arXiv: 1601.02077 [hep-th].\par
\noindent\hangindent=2.2em\hangafter=1 [84]\ O. Chacaltana, J. Distler, A. Trimm, and Y. Zhu. “Tinkertoys for the E7 theory.” JHEP 05 (2018), p. 031. doi: 10.1007/JHEP05(2018)031.arXiv: 1704.07890 [hep-th].\par
\noindent\hangindent=2.2em\hangafter=1 [85]\ O. Chacaltana, J. Distler, A. Trimm, and Y. Zhu. “Tinkertoys for the E8 Theory” (Feb. 2018). arXiv: 1802.09626 [hep-th].\par
\noindent\hangindent=2.2em\hangafter=1 [86]\ L. Charles and J. Marché. “Multicurves and regular functions on the representation variety of a surface in SU(2).” Comment. Math. Helv. 87.2 (2012), pp. 409--431. issn: 0010-2571. doi: 10.4171/CMH/258.arXiv: 0901.3064 [math.DG].\par
\noindent\hangindent=2.2em\hangafter=1 [87]\ M. Cho, S. Collier, and X. Yin. “Recursive Representations of Arbitrary Virasoro Conformal Blocks.” JHEP 04 (2019), p. 018. doi: 10.1007/JHEP04(2019)018.arXiv: 1703.09805 [hep-th].\par
\noindent\hangindent=2.2em\hangafter=1 [88]\ Y. Choi, C. Cordova, P.-S. Hsin, H. T. Lam, and S.-H. Shao. “Non-invertible Condensation, Duality, and Triality Defects in 3+1 Dimensions.” Commun. Math. Phys. 402.1 (2023), pp. 489--542. doi: 10.1007/s00220-023-04727-4.arXiv: 2204.09025 [hep-th].\par
\noindent\hangindent=2.2em\hangafter=1 [89]\ Y. Choi, H. T. Lam, and S.-H. Shao. “Noninvertible Global Symmetries in the Standard Model.” Phys. Rev. Lett. 129.16 (2022), p. 161601. doi: 10.1103/PhysRevLett.129.161601.arXiv: 2205.05086 [hep-th].\par
\noindent\hangindent=2.2em\hangafter=1 [90]\ C. Closset and H. Magureanu. “Reading between the rational sections: Global structures of 4d N = 2 KK theories.” SciPost Phys. 16.5 (2024), p. 137. doi: 10.21468/SciPostPhys.16.5.137.arXiv: 2308.10225 [hep-th].\par
\noindent\hangindent=2.2em\hangafter=1 [91]\ S. R. Coleman. “Classical Lumps and their Quantum Descendents.” Subnucl. Ser. 13 (1977), p. 297.\par
\noindent\hangindent=2.2em\hangafter=1 [92]\ S. R. Coleman. “THE MAGNETIC MONOPOLE FIFTY YEARS LATER.” Les Houches Summer School of Theoretical Physics: Laser-Plasma Interactions. June 1982, pp. 461-- 552.\par
\noindent\hangindent=2.2em\hangafter=1 [93]\ S. R. Coleman and D. J. Gross. “Price of asymptotic freedom.” Phys. Rev. Lett. 31 (1973), pp. 851--854. doi: 10.1103/PhysRevLett.31.851.\par
\noindent\hangindent=2.2em\hangafter=1 [94]\ S. R. Coleman and J. Mandula. “All Possible Symmetries of the S Matrix.” Phys. Rev. 159 (1967). Ed. by A. Zichichi, pp. 1251--1256. doi: 10.1103/PhysRev.159.1251.\par
\noindent\hangindent=2.2em\hangafter=1 [95]\ C. Córdova and S.-H. Shao. “Schur indices, BPS particles, and Argyres-Douglas theories.” J. High Energy Phys. 2016.1 (2016). Id/No 40, p. 38. issn: 1126-6708. doi: 10.1007/JHEP01(2016)040.\par
\noindent\hangindent=2.2em\hangafter=1 [96]\ K. Corlette. “Flat G-bundles with canonical metrics.” J. Differ. Geom. 28.3 (1988), pp. 361--382. issn: 0022-040X. doi: 10.4310/jdg/1214442469.\par
\noindent\hangindent=2.2em\hangafter=1 [97]\ C. W. Cryer. “The LU-factorization of totally positive matrices.” Linear Algebra Appl. 7 (1973), pp. 83--92. issn: 0024-3795. doi: 10.1016/0024-3795(73)90039-6.\par
\noindent\hangindent=2.2em\hangafter=1 [98]\ C. W. Cryer. “Some properties of totally positive matrices.” Linear Algebra Appl. 15 (1976), pp. 1--25. issn: 0024-3795. doi: 10.1016/0024-3795(76)90076-8.\par
\noindent\hangindent=2.2em\hangafter=1 [99]\ X.-z. Dai and D. S. Freed. “Eta invariants and determinant lines.” J. Math. Phys. 35 (1994). [Erratum: J.Math.Phys. 42, 2343--2344 (2001)], pp. 5155--5194. doi: 10.1063/1.530747.arXiv: hep-th/9405012.\par
\noindent\hangindent=2.2em\hangafter=1 [100]\ U. H. Danielsson and B. Sundborg. “The Moduli space and monodromies of N=2 supersymmetric SO(2r+1) Yang-Mills theory.” Phys. Lett. B 358 (1995), pp. 273--280. doi: 10.1016/0370-2693(95)01010-N.arXiv: hep-th/9504102.\par
\noindent\hangindent=2.2em\hangafter=1 [101]\ F. Del Monte and P. Longhi. “The threefold way to quantum periods: WKB, TBA equations and q-Painlevé.” SciPost Phys. 15.3 (2023), p. 112. doi: 10.21468/SciPostPhys.15.3.112.arXiv: 2207.07135 [hep-th].\par
\noindent\hangindent=2.2em\hangafter=1 [102]\ F. Del Monte and P. Longhi. “Monodromies of Second Order q-difference Equations from the WKB Approximation” (May 2024). arXiv: 2406.00175 [math-ph].\par
\noindent\hangindent=2.2em\hangafter=1 [103]\ P. Deligne and D. S. Freed. “Supersolutions.” Quantum fields and strings: a course for mathematicians. Vols. 1, 2. Material from the Special Year on Quantum Field Theory held at the Institute for Advanced Study, Princeton, NJ, 1996--1997. Providence, RI: AMS, American Mathematical Society, 1999, pp. 227--355. isbn: 0-8218-1987-9. arXiv: hep-th/9901094.\par
\noindent\hangindent=2.2em\hangafter=1 [104]\ M. Demazure. “Structures Algebriques --- Cohomologie Des Groupes.” Propriétés Générales des Schémas en Groupes. Berlin, Heidelberg: Springer Berlin Heidelberg, 1970, pp. 1--42. isbn: 978-3-540-36336-1.\par
\noindent\hangindent=2.2em\hangafter=1 [105]\ H. Derksen, J. Weyman, and A. Zelevinsky. “Quivers with potentials and their representations. I: Mutations.” Sel. Math., New Ser. 14.1 (2008), pp. 59--119. issn: 1022-1824. doi: 10.1007/s00029-008-0057-9.arXiv: 0704.0649 [math.RA].\par
\noindent\hangindent=2.2em\hangafter=1 [106]\ E. D’Hoker and D. H. Phong. “Lectures on supersymmetric Yang-Mills theory and integrable systems.” 9th CRM Summer School: Theoretical Physics at the End of the 20th Century. Dec. 1999, pp. 1--125. arXiv: hep-th/9912271.\par
\noindent\hangindent=2.2em\hangafter=1 [107]\ P. Di Francesco, C. Itzykson, and J. B. Zuber. “Classical W algebras.” Commun. Math. Phys. 140 (1991), pp. 543--568. doi: 10.1007/BF02099134.\par
\noindent\hangindent=2.2em\hangafter=1 [108]\ P. Di Francesco, P. Mathieu, and D. Senechal. Conformal Field Theory. Graduate Texts in Contemporary Physics. New York: Springer-Verlag, 1997. isbn: 978-1-4612-7475-9. doi: 10.1007/978-1-4612-2256-9.\par
\noindent\hangindent=2.2em\hangafter=1 [109]\ T. tom Dieck. Algebraic topology. EMS Textb. Math. Zürich: European Mathematical Society (EMS), 2008. isbn: 978-3-03719-048-7. doi: 10.4171/048.\par
\noindent\hangindent=2.2em\hangafter=1 [110]\ S. Dimopoulos and H. Georgi. “Softly Broken Supersymmetry and SU(5).” Nucl. Phys. B 193 (1981), pp. 150--162. doi: 10.1016/0550-3213(81)90522-8.\par
\noindent\hangindent=2.2em\hangafter=1 [111]\ M. Dine. Supersymmetry and String Theory: Beyond the Standard Model. Oxford University Press, 2016. isbn: 978-1-009-29089-0. doi: 10.1017/9781009290883.\par
\noindent\hangindent=2.2em\hangafter=1 [112]\ R. Donagi. “Decomposition of spectral covers.” Journées de géométrie algébrique d’Orsay, France, juillet 20-26, 1992. Paris: Société Mathématique de France, 1993, pp. 145--175.\par
\noindent\hangindent=2.2em\hangafter=1 [113]\ R. Donagi, L. Ein, and R. Lazarsfeld. “A Non-Linear Deformation of the Hitchin Dynamical System” (1995). arXiv: alg-geom/9504017.\par
\noindent\hangindent=2.2em\hangafter=1 [114]\ R. Donagi and E. Witten. “Supersymmetric Yang-Mills theory and integrable systems.” Nucl. Phys. B 460 (1996), pp. 299--334. doi: 10.1016/0550-3213(95)00609-5.arXiv: hep-th/9510101.\par
\noindent\hangindent=2.2em\hangafter=1 [115]\ S. K. Donaldson. “A new proof of a theorem of Narasimhan and Seshadri.” J. Differ. Geom. 18 (1983), pp. 269--277. issn: 0022-040X. doi: 10.4310/jdg/1214437664.\par
\noindent\hangindent=2.2em\hangafter=1 [116]\ S. K. Donaldson. “Twisted harmonic maps and the self-duality equations.” Proc. Lond. Math. Soc. (3) 55 (1987), pp. 127--131. issn: 0024-6115. doi: 10.1112/plms/s3-55.1.127.\par
\noindent\hangindent=2.2em\hangafter=1 [117]\ S. K. Donaldson and P. B. Kronheimer. The geometry of four-manifolds. Oxford Math. Monogr. Oxford: Clarendon Press, 1990. isbn: 0-19-853553-8.\par
\noindent\hangindent=2.2em\hangafter=1 [118]\ N. Dorey, T. J. Hollowood, V. V. Khoze, and M. P. Mattis. “The Calculus of many instantons.” Phys. Rept. 371 (2002), pp. 231--459. doi: 10.1016/S0370-1573(02)00301-0.arXiv: hep-th/0206063.\par
\noindent\hangindent=2.2em\hangafter=1 [119]\ P. Dorey, C. Dunning, and R. Tateo. “The ODE/IM Correspondence.” J. Phys. A 40 (2007), R205. doi: 10.1088/1751-8113/40/32/R01.arXiv: hep-th/0703066.\par
\noindent\hangindent=2.2em\hangafter=1 [120]\ P. Dorey and R. Tateo. “On the relation between Stokes multipliers and the T-Q systems of conformal field theory.” Nucl. Phys. B 563 (1999). [Erratum: Nucl.Phys.B 603, 581-- 581 (2001)], pp. 573--602. doi: 10.1016/S0550-3213(99)00609-4.arXiv: hepth/9906219.\par
\noindent\hangindent=2.2em\hangafter=1 [121]\ D. Dorigoni. “An Introduction to Resurgence, Trans-Series and Alien Calculus.” Annals Phys. 409 (2019), p. 167914. doi: 10.1016/j.aop.2019.167914.arXiv: 1411.3585 [hep-th].\par
\noindent\hangindent=2.2em\hangafter=1 [122]\ R. Eager, S. A. Selmani, and J. Walcher. “Exponential Networks and Representations of Quivers.” JHEP 08 (2017), p. 063. doi: 10.1007/JHEP08(2017)063.arXiv: 1611.06177 [hep-th].\par
\noindent\hangindent=2.2em\hangafter=1 [123]\ J. Écalle. Introduction aux fonctions analysables et preuve constructive de la conjecture de Dulac. Paris: Hermann, Éditeurs des Sciences et des Arts., 1992. isbn: 2-7056-6199-9.\par
\noindent\hangindent=2.2em\hangafter=1 [124]\ G. Faltings. “Stable G-bundles and projective connections.” J. Algebr. Geom. 2.3 (1993), pp. 507--568. issn: 1056-3911.\par
\noindent\hangindent=2.2em\hangafter=1 [125]\ B. Farb and D. Margalit. A primer on mapping class groups (pms-49). Vol. 41. Princeton university press, 2011.\par
\noindent\hangindent=2.2em\hangafter=1 [126]\ H. M. Farkas and I. Kra. Riemann surfaces. 2nd ed. Vol. 71. Grad. Texts Math. New York etc.: Springer-Verlag, 1992. isbn: 0-387-97703-1.\par
\noindent\hangindent=2.2em\hangafter=1 [127]\ F. Ferrari and A. Bilal. “The Strong coupling spectrum of the Seiberg-Witten theory.” Nucl. Phys. B 469 (1996), pp. 387--402. doi: 10.1016/0550-3213(96)00150-2.arXiv: hep-th/9602082.\par
\noindent\hangindent=2.2em\hangafter=1 [128]\ V. V. Fock. “Dual Teichmüller spaces” (1997). arXiv: dg-ga/9702018.\par
\noindent\hangindent=2.2em\hangafter=1 [129]\ V. Fock and A. Goncharov. “Moduli spaces of local systems and higher Teichmüller theory.” Publ. Math., Inst. Hautes Étud. Sci. 103 (2006), pp. 1--211. issn: 0073-8301. doi: 10.1007/s10240-006-0039-4.arXiv: math/0311149.\par
\noindent\hangindent=2.2em\hangafter=1 [130]\ V. Fock, V. Tatitscheff, and A. Thomas. “Topological quantum field theories from Hecke algebras.” Represent. Theory 27 (2023), pp. 248--291. issn: 1088-4165. doi: 10.1090/ert/640.arXiv: 2105.09622 [math.QA].\par
\noindent\hangindent=2.2em\hangafter=1 [131]\ V. V. Fock and A. B. Goncharov. “Dual Teichmüller and lamination spaces.” Handbook of Teichmüller theory. Volume I. Zürich: European Mathematical Society (EMS), 2007, pp. 647--684. isbn: 978-3-03719-029-6. arXiv: math/0510312.\par
\noindent\hangindent=2.2em\hangafter=1 [132]\ V. V. Fock and A. B. Goncharov. “Cluster ensembles, quantization and the dilogarithm.” Ann. Sci. Éc. Norm. Supér. (4) 42.6 (2009), pp. 865--930. issn: 0012-9593. doi: 10.24033/asens.2112.arXiv: math/0311245.\par
\noindent\hangindent=2.2em\hangafter=1 [133]\ S. Fomin and A. Zelevinsky. “Cluster algebras. I: Foundations.” J. Am. Math. Soc. 15.2 (2002), pp. 497--529. issn: 0894-0347. doi: 10.1090/S0894-0347-01-00385-X.arXiv: math/0104151.\par
\noindent\hangindent=2.2em\hangafter=1 [134]\ S. Fomin and A. Zelevinsky. “Cluster algebras. II: Finite type classification.” Invent. Math. 154.1 (2003), pp. 63--121. issn: 0020-9910. doi: 10.1007/s00222-003-0302-y.arXiv: math/0208229.\par
\noindent\hangindent=2.2em\hangafter=1 [135]\ S. Fomin and A. Zelevinsky. “Cluster algebras. IV: Coefficients.” Compos. Math. 143.1 (2007), pp. 112--164. issn: 0010-437X. doi: 10.1112/S0010437X06002521.arXiv: math/0602259.\par
\noindent\hangindent=2.2em\hangafter=1 [136]\ D. S. Freed. Five lectures on supersymmetry. Providence, USA: AMS, 1999. isbn: 978- 0-8218-1953-1.\par
\noindent\hangindent=2.2em\hangafter=1 [137]\ D. S. Freed. “Special Kahler manifolds.” Commun. Math. Phys. 203 (1999), pp. 31--52. doi: 10.1007/s002200050604.arXiv: hep-th/9712042.\par
\noindent\hangindent=2.2em\hangafter=1 [138]\ D. S. Freed and A. Neitzke. “3d spectral networks and classical Chern--Simons theory.” Surveys Diff. Geom. 26.1 (2021), pp. 51--155. doi: 10.4310/sdg.2021.v26.n1.a4.arXiv: 2208.07420 [math.DG].\par
\noindent\hangindent=2.2em\hangafter=1 [139]\ E. Frenkel and D. Ben-Zvi. Vertex algebras and algebraic curves. 2nd revised and expanded ed. Vol. 88. Math. Surv. Monogr. Providence, RI: American Mathematical Society (AMS), 2004. isbn: 0-8218-3674-9.\par
\noindent\hangindent=2.2em\hangafter=1 [140]\ R. Fricke and F. Klein. Vorlesungen über die Theorie der automorphen Funktionen. II. Bd.: Die funktionentheoretischen Ausführungen und die Anwendungen. 3. (Schluß) Lieferung. Direkte Beweismethoden der Fundamentaltheoreme und Anh. Leipzig: B. G. Teubner, 1912, pp. 439--668.\par
\noindent\hangindent=2.2em\hangafter=1 [141]\ K. Fujikawa. “Path Integral Measure for Gauge Invariant Fermion Theories.” Phys. Rev. Lett. 42 (1979), pp. 1195--1198. doi: 10.1103/PhysRevLett.42.1195.\par
\noindent\hangindent=2.2em\hangafter=1 [142]\ W. Fulton and J. Harris. Representation theory. A first course. Vol. 129. Grad. Texts Math. New York etc.: Springer-Verlag, 1991. isbn: 0-387-97527-6.\par
\noindent\hangindent=2.2em\hangafter=1 [143]\ M. Gabella. “BPS spectra from BPS graphs” (Oct. 2017). arXiv: 1710.08449 [hep-th].\par
\noindent\hangindent=2.2em\hangafter=1 [144]\ M. Gabella. “Quantum holonomies from spectral networks and framed BPS states.” Commun. Math. Phys. 351.2 (2017), pp. 563--598. issn: 0010-3616. doi: 10.1007/s00220-016-2729-1.\par
\noindent\hangindent=2.2em\hangafter=1 [145]\ M. Gabella, P. Longhi, C. Y. Park, and M. Yamazaki. “BPS Graphs: From Spectral Networks to BPS Quivers.” JHEP 07 (2017), p. 032. doi: 10.1007/JHEP07(2017)032.arXiv: 1704.04204 [hep-th].\par
\noindent\hangindent=2.2em\hangafter=1 [146]\ D. Gaiotto. “N=2 dualities.” JHEP 08 (2012), p. 034. doi: 10.1007/JHEP08(2012)034.arXiv: 0904.2715 [hep-th].\par
\noindent\hangindent=2.2em\hangafter=1 [147]\ D. Gaiotto. “Surface Operators in N = 2 4d Gauge Theories.” JHEP 11 (2012), p. 090. doi: 10.1007/JHEP11(2012)090.arXiv: 0911.1316 [hep-th].\par
\noindent\hangindent=2.2em\hangafter=1 [148]\ D. Gaiotto. “Opers and TBA” (Mar. 2014). arXiv: 1403.6137 [hep-th].\par
\noindent\hangindent=2.2em\hangafter=1 [149]\ D. Gaiotto, A. Kapustin, N. Seiberg, and B. Willett. “Generalized Global Symmetries.” JHEP 02 (2015), p. 172. doi: 10.1007/JHEP02(2015)172.arXiv: 1412. 5148 [hep-th].\par
\noindent\hangindent=2.2em\hangafter=1 [150]\ D. Gaiotto and J. Maldacena. “The Gravity duals of N=2 superconformal field theories.” JHEP 10 (2012), p. 189. doi: 10.1007/JHEP10(2012)189.arXiv: 0904.4466 [hep-th].\par
\noindent\hangindent=2.2em\hangafter=1 [151]\ D. Gaiotto, G. W. Moore, and A. Neitzke. “Four-dimensional wall-crossing via threedimensional field theory.” Commun. Math. Phys. 299 (2010), pp. 163--224. doi: 10.1007/s00220-010-1071-2.arXiv: 0807.4723 [hep-th].\par
\noindent\hangindent=2.2em\hangafter=1 [152]\ D. Gaiotto, G. W. Moore, and A. Neitzke. “Wall-Crossing in Coupled 2d-4d Systems.” JHEP 12 (2012), p. 082. doi: 10.1007/JHEP12(2012)082.arXiv: 1103. 2598 [hep-th].\par
\noindent\hangindent=2.2em\hangafter=1 [153]\ D. Gaiotto, G. W. Moore, and A. Neitzke. “Framed BPS States.” Adv. Theor. Math. Phys. 17.2 (2013), pp. 241--397. doi: 10.4310/ATMP.2013.v17.n2.a1.arXiv: 1006.0146 [hep-th].\par
\noindent\hangindent=2.2em\hangafter=1 [154]\ D. Gaiotto, G. W. Moore, and A. Neitzke. “Spectral networks.” Annales Henri Poincare 14 (2013), pp. 1643--1731. doi: 10.1007/s00023-013-0239-7.arXiv: 1204.4824 [hep-th].\par
\noindent\hangindent=2.2em\hangafter=1 [155]\ D. Gaiotto, G. W. Moore, and A. Neitzke. “Wall-crossing, Hitchin systems, and the WKB approximation.” Adv. Math. 234 (2013), pp. 239--403. doi: 10.1016/j.aim.2012.09.027.arXiv: 0907.3987 [hep-th].\par
\noindent\hangindent=2.2em\hangafter=1 [156]\ D. Gaiotto, G. W. Moore, and A. Neitzke. “Spectral Networks and Snakes.” Annales Henri Poincare 15 (2014), pp. 61--141. doi: 10.1007/s00023-013-0238-8.arXiv: 1209.0866 [hep-th].\par
\noindent\hangindent=2.2em\hangafter=1 [157]\ D. Gaiotto, G. W. Moore, A. Neitzke, and F. Yan. “Commuting line defects at qN= 1.” Commun. Math. Phys. 406.6 (2025). Id/No 144, p. 31. issn: 0010-3616. doi: 10.1007/s00220-025-05330-5.\par
\noindent\hangindent=2.2em\hangafter=1 [158]\ D. Gaiotto and E. Witten. “S-Duality of Boundary Conditions In N=4 Super Yang-Mills Theory.” Adv. Theor. Math. Phys. 13.3 (2009), pp. 721--896. doi: 10.4310/ATMP.2009.v13.n3.a5.arXiv: 0807.3720 [hep-th].\par
\noindent\hangindent=2.2em\hangafter=1 [159]\ I. n. García-Etxebarria and M. Montero. “Dai-Freed anomalies in particle physics.” JHEP 08 (2019), p. 003. doi: 10.1007/JHEP08(2019)003.arXiv: 1808.00009 [hep-th].\par
\noindent\hangindent=2.2em\hangafter=1 [160]\ I. n. García-Etxebarria and D. Regalado. “N = 3 four dimensional field theories.” JHEP 03 (2016), p. 083. doi: 10.1007/JHEP03(2016)083.arXiv: 1512.06434 [hep-th].\par
\noindent\hangindent=2.2em\hangafter=1 [161]\ I. n. García-Etxebarria and D. Regalado. “Exceptional N = 3 theories.” JHEP 12 (2017), p. 042. doi: 10.1007/JHEP12(2017)042.arXiv: 1611.05769 [hep-th].\par
\noindent\hangindent=2.2em\hangafter=1 [162]\ F. P. Gardiner. Teichmüller theory and quadratic differentials. Pure and Applied Mathematics. John Wiley \& Sons, 1987. isbn: 0-471-84539-6.\par
\noindent\hangindent=2.2em\hangafter=1 [163]\ M. Gasca and C. A. Micchelli, eds. Total positivity and its applications. Meeting, Zaragoza in Jaca, Spain, September 26--30, 1994. Vol. 359. Math. Appl., Dordr. Dordrecht: Kluwer Academic Publishers, 1996. isbn: 0-7923-3924-X.\par
\noindent\hangindent=2.2em\hangafter=1 [164]\ P. H. Ginsparg. “APPLIED CONFORMAL FIELD THEORY.” Les Houches Summer School in Theoretical Physics: Fields, Strings, Critical Phenomena. Sept. 1988. arXiv: hep-th/9108028.\par
\noindent\hangindent=2.2em\hangafter=1 [165]\ P. Goddard, J. Nuyts, and D. I. Olive. “Gauge Theories and Magnetic Charge.” Nucl. Phys. B 125 (1977), pp. 1--28. doi: 10.1016/0550-3213(77)90221-8.\par
\noindent\hangindent=2.2em\hangafter=1 [166]\ W. M. Goldman. “The symplectic nature of fundamental groups of surfaces.” Adv. Math. 54 (1984), pp. 200--225. issn: 0001-8708. doi: 10.1016/0001-8708(84)90040-9.\par
\noindent\hangindent=2.2em\hangafter=1 [167]\ W. M. Goldman. “Topological components of spaces of representations.” Invent. Math. 93.3 (1988), pp. 557--607. issn: 0020-9910. doi: 10.1007/BF01410200.\par
\noindent\hangindent=2.2em\hangafter=1 [168]\ A. Goncharov and M. Kontsevich. “Spectral description of non-commutative local systems on surfaces and non-commutative cluster varieties” (2021). arXiv: 2108.04168 [math.AG].\par
\noindent\hangindent=2.2em\hangafter=1 [169]\ A. Grassi, Q. Hao, and A. Neitzke. “Exponential Networks, WKB and Topological String.” SIGMA 19 (2023), p. 064. doi: 10.3842/SIGMA.2023.064.arXiv: 2201.11594 [hep-th].\par
\noindent\hangindent=2.2em\hangafter=1 [170]\ M. B. Green, J. H. Schwarz, and E. Witten. Superstring Theory Vol. 1: 25th Anniversary Edition. Cambridge Monographs on Mathematical Physics. Cambridge University Press, Nov. 2012. isbn: 978-1-107-02911-8. doi: 10.1017/CBO9781139248563.\par
\noindent\hangindent=2.2em\hangafter=1 [171]\ M. B. Green, J. H. Schwarz, and E. Witten. Superstring Theory Vol. 2: 25th Anniversary Edition. Cambridge Monographs on Mathematical Physics. Cambridge University Press, Nov. 2012. isbn: 978-1-107-02913-2. doi: 10.1017/CBO9781139248570.\par
\noindent\hangindent=2.2em\hangafter=1 [172]\ M. T. Grisaru, W. Siegel, and M. Rocek. “Improved Methods for Supergraphs.” Nucl. Phys. B 159 (1979), p. 429. doi: 10.1016/0550-3213(79)90344-4.\par
\noindent\hangindent=2.2em\hangafter=1 [173]\ D. J. Gross and F. Wilczek. “Asymptotically Free Gauge Theories - I.” Phys. Rev. D 8 (1973), pp. 3633--3652. doi: 10.1103/PhysRevD.8.3633.\par
\noindent\hangindent=2.2em\hangafter=1 [174]\ M. Gross, P. Hacking, S. Keel, and M. Kontsevich. “Canonical bases for cluster algebras.” J. Am. Math. Soc. 31.2 (2018), pp. 497--608. issn: 0894-0347. doi: 10.1090/jams/890.arXiv: 1411.1394 [math.AG].\par
\noindent\hangindent=2.2em\hangafter=1 [175]\ O. Guichard. “Composantes de Hitchin et représentations hyperconvexes de groupes de surface.” J. Differ. Geom. 80.3 (2008), pp. 391--431. issn: 0022-040X. doi: 10.4310/jdg/1226090482.\par
\noindent\hangindent=2.2em\hangafter=1 [176]\ O. Guichard. “An introduction to the differential geometry of flat bundles and of Higgs bundles.” The geometry, topology and physics of moduli spaces of Higgs bundles. Based on minicourses from the program on “The geometry, topology and physics of moduli spaces of Higgs bundles”, IMS, Singapore, July 7 -- August 29, 2014. Hackensack, NJ: World Scientific, 2018, pp. 1--63. isbn: 978-981-3229-10-5. doi: 10.1142/9789813229099\_0001.\par
\noindent\hangindent=2.2em\hangafter=1 [177]\ O. Guichard, F. Labourie, and A. Wienhard. “Positivity and representations of surface groups” (2021). arXiv: 2106.14584 [math.DG].\par
\noindent\hangindent=2.2em\hangafter=1 [178]\ O. Guichard, E. Rogozinnikov, and A. Wienhard. “Parametrizing spaces of positive representations” (2022). arXiv: 2210.11605 [math.DG].\par
\noindent\hangindent=2.2em\hangafter=1 [179]\ O. Guichard and A. Wienhard. “Positivity and higher Teichmüller theory.” European congress of mathematics. Proceedings of the 7th ECM (7ECM) congress, Berlin, Germany, July 18--22, 2016. Zürich: European Mathematical Society (EMS), 2018, pp. 289--310. isbn: 978-3-03719-676-2. doi: 10.4171/176-1/13.arXiv: 1802.02833 [math.DG].\par
\noindent\hangindent=2.2em\hangafter=1 [180]\ O. Guichard and A. Wienhard. Generalizing Lusztig’s total positivity. to appear in Inventiones Mathematicae. 2022. arXiv: 2208.10114 [math.DG].\par
\noindent\hangindent=2.2em\hangafter=1 [181]\ S. Gukov and E. Witten. “Gauge theory, ramification, and the geometric Langlands program.” Current developments in mathematics, 2006. Somerville, MA: International Press, 2008, pp. 35--180. isbn: 978-1-57146-167-4. arXiv: hep-th/0612073.\par
\noindent\hangindent=2.2em\hangafter=1 [182]\ R. Haag, J. T. Lopuszanski, and M. Sohnius. “All Possible Generators of Supersymmetries of the s Matrix.” Nucl. Phys. B 88 (1975), p. 257. doi: 10.1016/0550-3213(75)90279-5.\par
\noindent\hangindent=2.2em\hangafter=1 [183]\ F. Haiden, L. Katzarkov, and C. Simpson. “Spectral Networks and Stability Conditions for Fukaya Categories with Coefficients.” Commun. Math. Phys. 405.11 (2024), p. 264. doi: 10.1007/s00220-024-05138-9.arXiv: 2112.13623 [math.AG].\par
\noindent\hangindent=2.2em\hangafter=1 [184]\ A. Hanany. “On the quantum moduli space of vacua N=2 supersymmetric gauge theories.” Nucl. Phys. B 466 (1996), pp. 85--100. doi: 10.1016/0550-3213(96)00077-6.arXiv: hep-th/9509176.\par
\noindent\hangindent=2.2em\hangafter=1 [185]\ A. Hanany and Y. Oz. “On the quantum moduli space of vacua of N=2 supersymmetric SU(N(c)) gauge theories.” Nucl. Phys. B 452 (1995), pp. 283--312. doi: 10.1016/0550-3213(95)00376-4.arXiv: hep-th/9505075.\par
\noindent\hangindent=2.2em\hangafter=1 [186]\ A. Hanany and E. Witten. “Type IIB superstrings, BPS monopoles, and three-dimensional gauge dynamics.” Nucl. Phys. B 492 (1997), pp. 152--190. doi: 10.1016/S0550-3213(97)00157-0.arXiv: hep-th/9611230.\par
\noindent\hangindent=2.2em\hangafter=1 [187]\ Q. Hao, L. Hollands, and A. Neitzke. “BPS states in the Minahan-Nemeschansky E7 theory.” JHEP 04 (2020), p. 039. doi: 10.1007/JHEP04(2020)039.arXiv: 1905.09879 [hep-th].\par
\noindent\hangindent=2.2em\hangafter=1 [188]\ Q. Hao and A. Neitzke. “A new construction of c= 1 Virasoro blocks” (July 2024). arXiv: 2407.04483 [hep-th].\par
\noindent\hangindent=2.2em\hangafter=1 [189]\ R. Hartshorne. Algebraic geometry. Corr. 3rd printing. Vol. 52. Grad. Texts Math. Springer, Cham, 1983.\par
\noindent\hangindent=2.2em\hangafter=1 [190]\ J. A. Harvey. “Magnetic monopoles, duality and supersymmetry.” ICTP Summer School in High-energy Physics and Cosmology. Mar. 1996. arXiv: hep-th/9603086.\par
\noindent\hangindent=2.2em\hangafter=1 [191]\ J. A. Harvey. “TASI 2003 lectures on anomalies.” Sept. 2005. arXiv: hep-th/0509097.\par
\noindent\hangindent=2.2em\hangafter=1 [192]\ A. Hatcher. Algebraic topology. Cambridge: Cambridge University Press, 2002. isbn: 0-521-79540-0.\par
\noindent\hangindent=2.2em\hangafter=1 [193]\ A. E. Hatcher. “On triangulations of surfaces.” Topology Appl. 40.2 (1991), pp. 189--194. issn: 0166-8641. doi: 10.1016/0166-8641(91)90050-V.\par
\noindent\hangindent=2.2em\hangafter=1 [194]\ J. J. Heckman and T. Rudelius. “Top Down Approach to 6D SCFTs.” J. Phys. A 52.9 (2019), p. 093001. doi: 10.1088/1751-8121/aafc81.arXiv: 1805.06467 [hep-th].\par
\noindent\hangindent=2.2em\hangafter=1 [195]\ S. Helgason. Differential geometry, Lie groups, and symmetric spaces. Reprint with corrections of the 1978 original. Vol. 34. Grad. Stud. Math. Providence, RI: American Mathematical Society (AMS), 2001. isbn: 0-8218-2848-7.\par
\noindent\hangindent=2.2em\hangafter=1 [196]\ N. J. Hitchin, G. B. Segal, and R. S. Ward. Integrable systems. Twistors, loop groups, and Riemann surfaces. Based on lectures given at the instructional conference on integrable systems, Oxford, UK, September 1997. Reprint of the 1999 hardback ed. Vol. 4. Oxf. Grad. Texts Math. Oxford: Oxford University Press, 2013. isbn: 978-0-19-967677-4.\par
\noindent\hangindent=2.2em\hangafter=1 [197]\ N. Hitchin. “Stable bundles and integrable systems.” Duke Math. J. 54 (1987), pp. 91--114. issn: 0012-7094. doi: 10.1215/S0012-7094-87-05408-1.\par
\noindent\hangindent=2.2em\hangafter=1 [198]\ N. J. Hitchin. “The self-duality equations on a Riemann surface.” Proc. Lond. Math. Soc. (3) 55 (1987), pp. 59--126. issn: 0024-6115. doi: 10.1112/plms/s3-55.1.59.\par
\noindent\hangindent=2.2em\hangafter=1 [199]\ N. J. Hitchin. “Lie groups and Teichmüller space.” Topology 31.3 (1992), pp. 449--473. issn: 0040-9383. doi: 10.1016/0040-9383(92)90044-I.\par
\noindent\hangindent=2.2em\hangafter=1 [200]\ N. J. Hitchin. “The hyperkähler manifolds.” Séminaire Bourbaki, volume 1991/92. Exposés 745-759 (avec table par noms d’auteurs de 1948/49 à 1991/92). Paris: Société Mathématique de France, 1992, 137--166 (ex.\par
\noindent\hangindent=2.2em\hangafter=1 [201]\ N. J. Hitchin, A. Karlhede, U. Lindstrom, and M. Rocek. “Hyperkahler Metrics and Supersymmetry.” Commun. Math. Phys. 108 (1987), p. 535. doi: 10.1007/BF01214418.\par
\noindent\hangindent=2.2em\hangafter=1 [202]\ L. Hollands, C. A. Keller, and J. Song. “From SO/Sp instantons to W-algebra blocks.” J. High Energy Phys. 2011.3 (2011). Id/No 053, p. 80. issn: 1126-6708. doi: 10.1007/JHEP03(2011)053.\par
\noindent\hangindent=2.2em\hangafter=1 [203]\ L. Hollands, C. A. Keller, and J. Song. “Towards a 4d/2d correspondence for Sicilian quivers.” J. High Energy Phys. 2011.10 (2011). Id/No 100, p. 35. issn: 1126-6708. doi: 10.1007/JHEP10(2011)100.\par
\noindent\hangindent=2.2em\hangafter=1 [204]\ L. Hollands and O. Kidwai. “Higher length-twist coordinates, generalized Heun’s opers, and twisted superpotentials.” Adv. Theor. Math. Phys. 22.7 (2019), pp. 1713--1822. issn: 1095-0761. doi: 10.4310/ATMP.2018.v22.n7.a2.\par
\noindent\hangindent=2.2em\hangafter=1 [205]\ L. Hollands and A. Neitzke. “Spectral Networks and Fenchel--Nielsen Coordinates.” Lett. Math. Phys. 106.6 (2016), pp. 811--877. doi: 10.1007/s11005-016-0842-x.arXiv: 1312.2979 [math.GT].\par
\noindent\hangindent=2.2em\hangafter=1 [206]\ L. Hollands and A. Neitzke. “BPS states in the Minahan-Nemeschansky E6 theory.” Commun. Math. Phys. 353.1 (2017), pp. 317--351. doi: 10.1007/s00220-016-2798-1.arXiv: 1607.01743 [hep-th].\par
\noindent\hangindent=2.2em\hangafter=1 [207]\ L. Hollands and A. Neitzke. “Exact WKB and abelianization for the T3 equation.” Commun. Math. Phys. 380.1 (2020), pp. 131--186. doi: 10.1007/s00220-020-03875-1.arXiv: 1906.04271 [hep-th].\par
\noindent\hangindent=2.2em\hangafter=1 [208]\ L. Hollands, P. Rüter, and R. J. Szabo. “A geometric recipe for twisted superpotentials.” J. High Energy Phys. 2021.12 (2021). Id/No 164, p. 91. issn: 1126-6708. doi: 10.1007/JHEP12(2021)164.\par
\noindent\hangindent=2.2em\hangafter=1 [209]\ G. ’t Hooft. “Magnetic Monopoles in Unified Gauge Theories.” Nucl. Phys. B 79 (1974). Ed. by J. C. Taylor, pp. 276--284. doi: 10.1016/0550-3213(74)90486-6.\par
\noindent\hangindent=2.2em\hangafter=1 [210]\ J. Horn and M. Möller. Wall-crossing formulas via spectral networks. 2025. arXiv: 2508.07727 [math.AG].\par
\noindent\hangindent=2.2em\hangafter=1 [211]\ J. H. Hubbard. Teichmüller theory and applications to geometry, topology, and dynamics. Volume 1--3. Ithaca, NY: Matrix Editions, 2022.\par
\noindent\hangindent=2.2em\hangafter=1 [212]\ Y. Imayoshi and M. Taniguchi. An introduction to Teichmüller spaces. Tokyo: Springer- Verlag, 1992. isbn: 4-431-70088-9.\par
\noindent\hangindent=2.2em\hangafter=1 [213]\ K. A. Intriligator and N. Seiberg. “Lectures on supersymmetric gauge theories and electricmagnetic duality.” Nucl. Phys. B Proc. Suppl. 45BC (1996). Ed. by A. Zichichi, pp. 1--28. doi: 10.1016/0920-5632(95)00626-5.arXiv: hep-th/9509066.\par
\noindent\hangindent=2.2em\hangafter=1 [214]\ M. Ionita. “Spectral Networks and non-abelianization for reductive groups.” PhD thesis. University of Pennsylvania, 2020.\par
\noindent\hangindent=2.2em\hangafter=1 [215]\ M. Ionita and B. Morrissey. “Spectral Networks and Non-abelianization” (Mar. 2021). arXiv: 2103.12285 [math.AG].\par
\noindent\hangindent=2.2em\hangafter=1 [216]\ K. Iwaki and T. Nakanishi. “Exact WKB analysis and cluster algebras.” J. Phys. A 47.47 (2014), p. 474009. doi: 10.1088/1751-8113/47/47/474009.arXiv: 1401.7094 [math.CA].\par
\noindent\hangindent=2.2em\hangafter=1 [217]\ C. V. Johnson. D-Branes. Cambridge Monographs on Mathematical Physics. 2023. isbn: 978-0-511-60654-0. doi: 10.1017/9781009401371.\par
\noindent\hangindent=2.2em\hangafter=1 [218]\ J. Jost. Compact Riemann surfaces. An introduction to contemporary mathematics. 3rd ed. Universitext. Berlin: Springer, 2006. isbn: 3-540-33065-8. doi: 10.1007/978-3-540-33067-7.\par
\noindent\hangindent=2.2em\hangafter=1 [219]\ B. Julia and A. Zee. “Poles with Both Magnetic and Electric Charges in Nonabelian Gauge Theory.” Phys. Rev. D 11 (1975), pp. 2227--2232. doi: 10.1103/PhysRevD.11.2227.\par
\noindent\hangindent=2.2em\hangafter=1 [220]\ J. Kaidi, K. Ohmori, and Y. Zheng. “Kramers-Wannier-like Duality Defects in (3+1)D Gauge Theories.” Phys. Rev. Lett. 128.11 (2022), p. 111601. doi: 10.1103/PhysRevLett.128.111601.arXiv: 2111.01141 [hep-th].\par
\noindent\hangindent=2.2em\hangafter=1 [221]\ J. Kaidi, G. Zafrir, and Y. Zheng. “Non-invertible symmetries of N = 4 SYM and twisted compactification.” JHEP 08 (2022), p. 053. doi: 10.1007/JHEP08(2022)053.arXiv: 2205.01104 [hep-th].\par
\noindent\hangindent=2.2em\hangafter=1 [222]\ I. Kapovich and N. Benakli. “Boundaries of hyperbolic groups.” Combinatorial and geometric group theory. Proceedings of the AMS special session on combinatorial group theory, New York, NY, USA, November 4--5, 2000 and the AMS special session on computational group theory, Hoboken, NJ, USA, April 28--29, 2001. Providence, RI: American Mathematical Society (AMS), 2002, pp. 39--93. isbn: 0-8218-2822-3.\par
\noindent\hangindent=2.2em\hangafter=1 [223]\ M. Kapovich, B. Leeb, and J. Porti. “A Morse lemma for quasigeodesics in symmetric spaces and Euclidean buildings.” Geom. Topol. 22.7 (2018), pp. 3827--3923. issn: 1465- 3060. doi: 10.2140/gt.2018.22.3827.\par
\noindent\hangindent=2.2em\hangafter=1 [224]\ A. Kapustin. “D(n) quivers from branes.” JHEP 12 (1998), p. 015. doi: 10.1088/1126-6708/1998/12/015.arXiv: hep-th/9806238.\par
\noindent\hangindent=2.2em\hangafter=1 [225]\ A. Kapustin. “Wilson-’t Hooft operators in four-dimensional gauge theories and Sduality.” Phys. Rev. D 74 (2006), p. 025005. doi: 10.1103/PhysRevD.74.025005.arXiv: hep-th/0501015.\par
\noindent\hangindent=2.2em\hangafter=1 [226]\ A. Kapustin and E. Witten. “Electric-Magnetic Duality And The Geometric Langlands Program.” Commun. Num. Theor. Phys. 1 (2007), pp. 1--236. doi: 10.4310/CNTP.2007.v1.n1.a1.arXiv: hep-th/0604151.\par
\noindent\hangindent=2.2em\hangafter=1 [227]\ S. Karlin. Total positivity. Volume I. Stanford University Press, 1968.\par
\noindent\hangindent=2.2em\hangafter=1 [228]\ R. M. Kashaev. “Quantization of Teichmüller spaces and the quantum dilogarithm.” Lett. Math. Phys. 43.2 (1998), pp. 105--115. issn: 0377-9017. doi: 10.1023/A:1007460128279.\par
\noindent\hangindent=2.2em\hangafter=1 [229]\ T. Kawai and Y. Takei. Algebraic analysis of singular perturbation theory. Transl. from the Japanese by Goro Kato. Vol. 227. Transl. Math. Monogr. Providence, RI: American Mathematical Society (AMS), 2005. isbn: 0-8218-3547-5.\par
\noindent\hangindent=2.2em\hangafter=1 [230]\ C. A. Keller, N. Mekareeya, J. Song, and Y. Tachikawa. “The ABCDEFG of Instantons and W-algebras.” JHEP 03 (2012), p. 045. doi: 10.1007/JHEP03(2012)045.arXiv: 1111.5624 [hep-th].\par
\noindent\hangindent=2.2em\hangafter=1 [231]\ H. K. Kim, T. T. Q. Lê, and M. Son. “SL2 quantum trace in quantum Teichmüller theory via writhe.” Algebr. Geom. Topol. 23.1 (2023), pp. 339--418. issn: 1472-2747. doi: 10.2140/agt.2023.23.339.\par
\noindent\hangindent=2.2em\hangafter=1 [232]\ A. Klemm, W. Lerche, S. Yankielowicz, and S. Theisen. “Simple singularities and N=2 supersymmetric Yang-Mills theory.” Phys. Lett. B 344 (1995), pp. 169--175. doi: 10.1016/0370-2693(94)01516-F.arXiv: hep-th/9411048.\par
\noindent\hangindent=2.2em\hangafter=1 [233]\ A. Klemm, W. Lerche, P. Mayr, C. Vafa, and N. P. Warner. “Selfdual strings and N=2 supersymmetric field theory.” Nucl. Phys. B 477 (1996), pp. 746--766. doi: 10.1016/0550-3213(96)00353-7.arXiv: hep-th/9604034.\par
\noindent\hangindent=2.2em\hangafter=1 [234]\ A. W. Knapp. Lie groups beyond an introduction. 2nd ed. Vol. 140. Prog. Math. Boston, MA: Birkhäuser, 2002. isbn: 0-8176-4259-5.\par
\noindent\hangindent=2.2em\hangafter=1 [235]\ S. Kobayashi and K. Nomizu. Foundations of differential geometry. I. Vol. 15. Intersci. Tracts Pure Appl. Math. Interscience Publishers, New York, NY, 1963.\par
\noindent\hangindent=2.2em\hangafter=1 [236]\ K. Kodaira. “On the structure of compact complex analytic surfaces. II, III.” Am. J. Math. 88 (1966), pp. 682--721. issn: 0002-9327. doi: 10.2307/2373150.\par
\noindent\hangindent=2.2em\hangafter=1 [237]\ M. Kontsevich and Y. Soibelman. “Stability structures, motivic Donaldson-Thomas invariants and cluster transformations” (Nov. 2008). arXiv: 0811.2435 [math.AG].\par
\noindent\hangindent=2.2em\hangafter=1 [238]\ J. Korinman and A. Quesney. “The quantum trace as a quantum non-abelianization map.” J. Knot Theory Ramifications 31.6 (2022). Id/No 2250032, p. 49. issn: 0218-2165. doi: 10.1142/S0218216522500328.\par
\noindent\hangindent=2.2em\hangafter=1 [239]\ B. Kostant. “The principal three-dimensional subgroup and the Betti numbers of a complex simple Lie group.” Am. J. Math. 81 (1959), pp. 973--1032. issn: 0002-9327. doi: 10.2307/2372999.\par
\noindent\hangindent=2.2em\hangafter=1 [240]\ B. Kostant. “Root systems for Levi factors and Borel-de Siebenthal theory.” Symmetry and spaces. In Honor of Gerry Schwarz on the occasion of his 60th birthday. Basel: Birkhäuser, 2010, pp. 129--152. isbn: 978-0-8176-4875-6. doi: 10.1007/978-0-8176-4875-6\_7.\par
\noindent\hangindent=2.2em\hangafter=1 [241]\ H. A. Kramers. “Wellenmechanik und halbzahlige Quantisierung.” Z. Phys. 39 (1926), pp. 828--840. issn: 0939-7922.\par
\noindent\hangindent=2.2em\hangafter=1 [242]\ F. Labourie. “Anosov flows, surface groups and curves in projective space.” Invent. Math. 165.1 (2006), pp. 51--114. issn: 0020-9910. doi: 10.1007/s00222-005-0487-3.\par
\noindent\hangindent=2.2em\hangafter=1 [243]\ I. Le. “Higher laminations and affine buildings.” Geom. Topol. 20.3 (2016), pp. 1673-- 1735. issn: 1465-3060. doi: 10.2140/gt.2016.20.1673.arXiv: 1209. 0812 [math.RT].\par
\noindent\hangindent=2.2em\hangafter=1 [244]\ B. Le Floch. “A slow review of the AGT correspondence.” J. Phys. A 55.35 (2022), p. 353002. doi: 10.1088/1751-8121/ac5945.arXiv: 2006.14025 [hep-th].\par
\noindent\hangindent=2.2em\hangafter=1 [245]\ W. Lerche. “Introduction to Seiberg-Witten theory and its stringy origin.” AIP Conf. Proc. 419.1 (1998). Ed. by H. Falomir, R. E. Gamboa Saravi, and F. A. Schaposnik, pp. 171--217. doi: 10.1063/1.54690.\par
\noindent\hangindent=2.2em\hangafter=1 [246]\ A. Lesfari. “Adjunction formula, Poincaré residue and holomorphic differentials on Riemann surfaces.” Methods Funct. Anal. Topol. 24.1 (2018), pp. 41--52. issn: 1029-3531.\par
\noindent\hangindent=2.2em\hangafter=1 [247]\ J. Li. “The space of surface group representations.” Manuscr. Math. 78.3 (1993), pp. 223-- 243. issn: 0025-2611. doi: 10.1007/BF02599310.\par
\noindent\hangindent=2.2em\hangafter=1 [248]\ C. Loewner. “On totally positive matrices.” Math. Z. 63 (1955), pp. 338--340. issn: 0025- 5874. doi: 10.1007/BF01187945.\par
\noindent\hangindent=2.2em\hangafter=1 [249]\ P. Longhi and C. Y. Park. “ADE Spectral Networks.” JHEP 08 (2016), p. 087. doi: 10.1007/JHEP08(2016)087.arXiv: 1601.02633 [hep-th].\par
\noindent\hangindent=2.2em\hangafter=1 [250]\ P. Longhi, C. Y. Park, and A. Neitzke. loom. https://github.com/chan-y-park/ loom. 2016.\par
\noindent\hangindent=2.2em\hangafter=1 [251]\ G. Lusztig. “Total positivity in reductive groups.” Lie theory and geometry: in honor of Bertram Kostant on the occasion of his 65th birthday. Invited papers, some originated at a symposium held at MIT, Cambridge, MA, USA in May 1993. Boston, MA: Birkhäuser, 1994, pp. 531--568. isbn: 0-8176-3761-3.\par
\noindent\hangindent=2.2em\hangafter=1 [252]\ M. A. Luty and W. Taylor. “Varieties of vacua in classical supersymmetric gauge theories.” Phys. Rev. D 53 (1996), pp. 3399--3405. doi: 10.1103/PhysRevD.53.3399.arXiv: hep-th/9506098.\par
\noindent\hangindent=2.2em\hangafter=1 [253]\ J. D. Lykken. “Introduction to supersymmetry.” Theoretical Advanced Study Institute in Elementary Particle Physics (TASI 96): Fields, Strings, and Duality. June 1996, pp. 85-- 153. arXiv: hep-th/9612114.\par
\noindent\hangindent=2.2em\hangafter=1 [254]\ J. M. Maldacena. “The Large N limit of superconformal field theories and supergravity.” Adv. Theor. Math. Phys. 2 (1998), pp. 231--252. doi: 10.4310/ATMP.1998.v2.n2.a1.arXiv: hep-th/9711200.\par
\noindent\hangindent=2.2em\hangafter=1 [255]\ J. Marché and M. Narimannejad. “Some asymptotics of topological quantum field theory via skein theory.” Duke Math. J. 141.3 (2008), pp. 573--587. issn: 0012-7094. doi: 10.1215/00127094-2007-006.arXiv: math/0604533.\par
\noindent\hangindent=2.2em\hangafter=1 [256]\ A. Marden. Hyperbolic manifolds. An introduction in 2 and 3 dimensions. 2nd edition of the book previously published under the title ‘Outer circles. An introduction to hyperbolic 3-manifolds’. Cambridge: Cambridge University Press, 2016. isbn: 978-1-316-33777-6. doi: 10.1017/CBO9781316337776.\par
\noindent\hangindent=2.2em\hangafter=1 [257]\ M. Mariño. “Lectures on non-perturbative effects in large Ngauge theories, matrix models and strings.” Fortsch. Phys. 62 (2014), pp. 455--540. doi: 10.1002/prop.201400005.arXiv: 1206.6272 [hep-th].\par
\noindent\hangindent=2.2em\hangafter=1 [258]\ M. Mariño. Advanced topics in quantum mechanics. Cambridge: Cambridge University Press, 2022. isbn: 978-1-108-49587-5. doi: 10.1017/9781108863384.\par
\noindent\hangindent=2.2em\hangafter=1 [259]\ E. Markman. “Spectral curves and integrable systems.” Compos. Math. 93.3 (1994), pp. 255--290. issn: 0010-437X.\par
\noindent\hangindent=2.2em\hangafter=1 [260]\ R. J. Marsh. Lecture notes on cluster algebras. Zur. Lect. Adv. Math. Zürich: European Mathematical Society (EMS), 2013. isbn: 978-3-03719-130-9. doi: 10.4171/130.\par
\noindent\hangindent=2.2em\hangafter=1 [261]\ B. Martelli. An Introduction to Geometric Topology. 2016. arXiv: 1610. 02592 [math.GT].\par
\noindent\hangindent=2.2em\hangafter=1 [262]\ S. P. Martin. “A Supersymmetry primer.” Adv. Ser. Direct. High Energy Phys. 18 (1998). Ed. by G. L. Kane, pp. 1--98. doi: 10.1142/9789812839657\_0001.arXiv: hepph/9709356.\par
\noindent\hangindent=2.2em\hangafter=1 [263]\ K. Maruyoshi, C. Y. Park, and W. Yan. “BPS spectrum of Argyres-Douglas theory via spectral network.” JHEP 12 (2013), p. 092. doi: 10.1007/JHEP12(2013)092.arXiv: 1309.3050 [hep-th].\par
\noindent\hangindent=2.2em\hangafter=1 [264]\ W. S. Massey. A basic course in algebraic topology. Vol. 127. Grad. Texts Math. New York etc.: Springer-Verlag, 1991. isbn: 0-387-97430-X.\par
\noindent\hangindent=2.2em\hangafter=1 [265]\ A. Mikhailov. “BPS states and minimal surfaces.” Nucl. Phys. B 533 (1998), pp. 243--274. doi: 10.1016/S0550-3213(98)00524-0.arXiv: hep-th/9708068.\par
\noindent\hangindent=2.2em\hangafter=1 [266]\ J. A. Minahan and D. Nemeschansky. “An N=2 superconformal fixed point with E(6) global symmetry.” Nucl. Phys. B 482 (1996), pp. 142--152. doi: 10.1016/S0550-3213(96)00552-4.arXiv: hep-th/9608047.\par
\noindent\hangindent=2.2em\hangafter=1 [267]\ J. A. Minahan and D. Nemeschansky. “Superconformal fixed points with E(n) global symmetry.” Nucl. Phys. B 489 (1997), pp. 24--46. doi: 10.1016/S0550-3213(97)00039-4.arXiv: hep-th/9610076.\par
\noindent\hangindent=2.2em\hangafter=1 [268]\ S. Minwalla. “Restrictions imposed by superconformal invariance on quantum field theories.” Adv. Theor. Math. Phys. 2 (1998), pp. 783--851. doi: 10.4310/ATMP.1998.v2.n4.a4.arXiv: hep-th/9712074.\par
\noindent\hangindent=2.2em\hangafter=1 [269]\ R. Miranda. Algebraic curves and Riemann surfaces. Vol. 5. Grad. Stud. Math. Providence, RI: AMS, American Mathematical Society, 1995. isbn: 0-8218-0268-2.\par
\noindent\hangindent=2.2em\hangafter=1 [270]\ S. Monnier. “A Modern Point of View on Anomalies.” Fortsch. Phys. 67.8-9 (2019), p. 1910012. doi: 10.1002/prop.201910012.arXiv: 1903.02828 [hep-th].\par
\noindent\hangindent=2.2em\hangafter=1 [271]\ C. Montonen and D. I. Olive. “Magnetic Monopoles as Gauge Particles?” Phys. Lett. B 72 (1977), pp. 117--120. doi: 10.1016/0370-2693(77)90076-4.\par
\noindent\hangindent=2.2em\hangafter=1 [272]\ G. W. Moore and Y. Tachikawa. “On 2d TQFTs whose values are holomorphic symplectic varieties.” Proc. Symp. Pure Math. 85 (2012). Ed. by J. Block, J. Distler, R. Donagi, and E. Sharpe, pp. 191--208. doi: 10.1090/pspum/085/1379.arXiv: 1106.5698 [hep-th].\par
\noindent\hangindent=2.2em\hangafter=1 [273]\ J. W. Morgan. “An introduction to gauge theory.” Gauge theory and the topology of four-manifolds. Lectures of the graduate summer school. Providence, RI: American Mathematical Society, 1998, pp. 53--143. isbn: 0-8218-0591-6.\par
\noindent\hangindent=2.2em\hangafter=1 [274]\ B. J. R. Morrissey. “Nonabelianization, Spectral data and cameral data.” PhD thesis. University of Pennsylvania, 2020.\par
\noindent\hangindent=2.2em\hangafter=1 [275]\ L. Mosher. “Tiling the projective foliation space of a punctured surface.” Trans. Am. Math. Soc. 306.1 (1988), pp. 1--70. issn: 0002-9947. doi: 10.2307/2000830.\par
\noindent\hangindent=2.2em\hangafter=1 [276]\ D. Mumford. “Projective invariants of projective structures and applications.” Proc. Int. Congr. Math. 1962. 1963, pp. 526--530.\par
\noindent\hangindent=2.2em\hangafter=1 [277]\ J. R. Munkres. Elements of algebraic topology. Advanced Book Program. Addison-Wesley Publishing Company, Inc. (The Benjamin/Cummings Publ. Co., Inc.), 1984. isbn: 0-201- 04586-9.\par
\noindent\hangindent=2.2em\hangafter=1 [278]\ M. S. Narasimhan and C. S. Seshadri. “Stable and unitary vector bundles on a compact Riemann surface.” Ann. Math. (2) 82 (1965), pp. 540--567. issn: 0003-486X. doi: 10.2307/1970710.\par
\noindent\hangindent=2.2em\hangafter=1 [279]\ A. Neitzke. “Lectures on BPS states and spectral networks.” Quantum field theory and manifold invariants. Lecture notes from the IAS/Park City Mathematics Institute (PCMI) summer school, Park City, UT, USA, 2019. Providence, RI: American Mathematical Society (AMS), 2021, pp. 439--476. isbn: 978-1-4704-6123-2. doi: 10.1090/pcms/028/00881.\par
\noindent\hangindent=2.2em\hangafter=1 [280]\ A. Neitzke. “Integral iterations for harmonic maps.” Beijing J. Pure Appl. Math. 1.1 (2024), pp. 121--158. doi: 10.4310/bpam.2024.v1.n1.a3.arXiv: 1704.01522 [math.DG].\par
\noindent\hangindent=2.2em\hangafter=1 [281]\ A. Neitzke and F. Yan. “q-nonabelianization for line defects.” J. High Energy Phys. 2020.9 (2020). Id/No 153, p. 66. issn: 1126-6708. doi: 10.1007/JHEP09(2020)153.\par
\noindent\hangindent=2.2em\hangafter=1 [282]\ A. Neitzke and F. Yan. “The quantum UV-IR map for line defects in gl(3)-type class S theories.” J. High Energy Phys. 2022.9 (2022). Id/No 81, p. 51. issn: 1126-6708. doi: 10.1007/JHEP09(2022)081.\par
\noindent\hangindent=2.2em\hangafter=1 [283]\ N. Nekrasov. “Five dimensional gauge theories and relativistic integrable systems.” Nucl. Phys. B 531 (1998), pp. 323--344. doi: 10.1016/S0550-3213(98)00436-2.arXiv: hep-th/9609219.\par
\noindent\hangindent=2.2em\hangafter=1 [284]\ N. A. Nekrasov. “Seiberg-Witten prepotential from instanton counting.” Adv. Theor. Math. Phys. 7.5 (2003), pp. 831--864. doi: 10.4310/ATMP.2003.v7.n5.a4.arXiv: hepth/0206161.\par
\noindent\hangindent=2.2em\hangafter=1 [285]\ G. Nemes. “On the Borel summability of WKB solutions of certain Schrödinger-type differential equations.” J. Approx. Theory 265 (2021). Id/No 105562, p. 30. issn: 0021- 9045. doi: 10.1016/j.jat.2021.105562.arXiv: 2004.13367 [math.CA].\par
\noindent\hangindent=2.2em\hangafter=1 [286]\ B. C. Ngô. “Le lemme fondamental pour les algèbres de Lie.” Publ. Math., Inst. Hautes Étud. Sci. 111 (2010), pp. 1--169. issn: 0073-8301. doi: 10.1007/s10240-010-0026-7.arXiv: 0801.0446 [math.AG].\par
\noindent\hangindent=2.2em\hangafter=1 [287]\ Y. J. Nho. Family Floer theory, non-abelianization, and Spectral Networks. Preprint, arXiv: 2307.04213 [math.SG] (2023). 2023.\par
\noindent\hangindent=2.2em\hangafter=1 [288]\ N. Nikolaev. “Abelianisation of logarithmic sl2-connections.” Sel. Math., New Ser. 27.5 (2021). Id/No 78, p. 35. issn: 1022-1824. doi: 10.1007/s00029-021-00688-5.\par
\noindent\hangindent=2.2em\hangafter=1 [289]\ N. Nikolaev. “Existence and Uniqueness of Exact WKB Solutions for Second-Order Singularly Perturbed Linear ODEs.” Commun. Math. Phys. 400.1 (2023), pp. 463--517. doi: 10.1007/s00220-022-04603-7.arXiv: 2106.10248 [math.AP].\par
\noindent\hangindent=2.2em\hangafter=1 [290]\ N. Nikolaev. “Geometry and Resurgence of WKB Solutions of Schrödinger Equations” (2024).\par
\noindent\hangindent=2.2em\hangafter=1 [291]\ A. L. Onishchik. Lectures on real semisimple Lie algebras and their representations. ESI Lect. Math. Phys. Zürich: European Mathematical Society Publishing House, 2004. isbn: 3-03719-002-7. doi: 10.4171/002.\par
\noindent\hangindent=2.2em\hangafter=1 [292]\ H. Osborn. “Topological Charges for N=4 Supersymmetric Gauge Theories and Monopoles of Spin 1.” Phys. Lett. B 83 (1979), pp. 321--326. doi: 10.1016/0370-2693(79)91118-3.\par
\noindent\hangindent=2.2em\hangafter=1 [293]\ V. Ovsienko and S. Tabachnikov. “What is... the Schwarzian derivative?” Notices Am. Math. Soc. 56.1 (2009), pp. 34--36. issn: 0002-9920.\par
\noindent\hangindent=2.2em\hangafter=1 [294]\ T. Pantev and E. Sharpe. “Duality group actions on fermions.” JHEP 11 (2016), p. 171. doi: 10.1007/JHEP11(2016)171.arXiv: 1609.00011 [hep-th].\par
\noindent\hangindent=2.2em\hangafter=1 [295]\ A. Papadopoulos, ed. Handbook of Teichmüller theory. Volume I--VII. Vol. 11--30. IRMA Lect. Math. Theor. Phys. Berlin: European Mathematical Society (EMS), 2020.\par
\noindent\hangindent=2.2em\hangafter=1 [296]\ A. Papadopoulos. “Ideal triangles, hyperbolic surfaces and the Thurston metric on Teichmüller space.” Moduli spaces and locally symmetric spaces. Based on two workshops, Morningside Center of Mathematics, Beijing, China, February 2017 and March 2019. Somerville, MA: International Press; Beijing: Higher Education Press, 2021, pp. 135-- 181. isbn: 978-1-57146-409-5. arXiv: 2103.10066 [math.GT].\par
\noindent\hangindent=2.2em\hangafter=1 [297]\ R. C. Penner. “The decorated Teichmüller space of punctured surfaces.” Commun. Math. Phys. 113 (1987), pp. 299--339. issn: 0010-3616. doi: 10.1007/BF01223515.\par
\noindent\hangindent=2.2em\hangafter=1 [298]\ R. C. Penner. Decorated Teichmüller theory. QGM Master Cl. Ser. Zürich: European Mathematical Society (EMS), 2012. isbn: 978-3-03719-075-3. doi: 10.4171/075.\par
\noindent\hangindent=2.2em\hangafter=1 [299]\ R. Penrose. “The twistor programme.” Reports on Mathematical Physics 12.1 (1977), pp. 65--76. issn: 0034-4877. doi: https://doi.org/10.1016/0034-4877(77)90047-7.\par
\noindent\hangindent=2.2em\hangafter=1 [300]\ M. E. Peskin. “Duality in supersymmetric Yang-Mills theory.” Theoretical Advanced Study Institute in Elementary Particle Physics (TASI 96): Fields, Strings, and Duality. Feb. 1997, pp. 729--809. arXiv: hep-th/9702094.\par
\noindent\hangindent=2.2em\hangafter=1 [301]\ J. Polchinski. String theory. Vol. 1: An introduction to the bosonic string. Cambridge Monographs on Mathematical Physics. Cambridge University Press, Dec. 2007. isbn: 978-0-521-63303-1. doi: 10.1017/CBO9780511816079.\par
\noindent\hangindent=2.2em\hangafter=1 [302]\ J. Polchinski. String theory. Vol. 2: Superstring theory and beyond. Cambridge Monographs on Mathematical Physics. Cambridge University Press, Dec. 2007. isbn: 978-0- 521-67228-3. doi: 10.1017/CBO9780511618123.\par
\noindent\hangindent=2.2em\hangafter=1 [303]\ J. Polchinski. “Dirichlet Branes and Ramond-Ramond charges.” Phys. Rev. Lett. 75 (1995), pp. 4724--4727. doi: 10.1103/PhysRevLett.75.4724.arXiv: hep-th/9510017.\par
\noindent\hangindent=2.2em\hangafter=1 [304]\ H. D. Politzer. “Reliable Perturbative Results for Strong Interactions?” Phys. Rev. Lett. 30 (1973). Ed. by J. C. Taylor, pp. 1346--1349. doi: 10.1103/PhysRevLett.30.1346.\par
\noindent\hangindent=2.2em\hangafter=1 [305]\ A. M. Polyakov. “Particle Spectrum in Quantum Field Theory.” JETP Lett. 20 (1974). Ed. by J. C. Taylor, pp. 194--195.\par
\noindent\hangindent=2.2em\hangafter=1 [306]\ M. K. Prasad and C. M. Sommerfield. “An Exact Classical Solution for the ’t Hooft Monopole and the Julia-Zee Dyon.” Phys. Rev. Lett. 35 (1975), pp. 760--762. doi: 10.1103/PhysRevLett.35.760.\par
\noindent\hangindent=2.2em\hangafter=1 [307]\ A. Ramanathan. “Stable principal bundles on a compact Riemann surface.” Math. Ann. 213 (1975), pp. 129--152. issn: 0025-5831. doi: 10.1007/BF01343949.\par
\noindent\hangindent=2.2em\hangafter=1 [308]\ A. Ramsay and R. D. Richtmyer. Introduction to hyperbolic geometry. Universitext. New York, NY: Springer-Verlag, 1995. isbn: 0-387-94339-0.\par
\noindent\hangindent=2.2em\hangafter=1 [309]\ J. G. Ratcliffe. Foundations of hyperbolic manifolds. 3rd expanded edition. Vol. 149. Grad. Texts Math. Cham: Springer, 2019. isbn: 978-3-030-31597-9. doi: 10.1007/978-3-030-31597-9.\par
\noindent\hangindent=2.2em\hangafter=1 [310]\ D. Sauzin. “Introduction to 1---Summability and Resurgence.” Divergent Series, Summability and Resurgence I: Monodromy and Resurgence. Cham: Springer International Publishing, 2016, pp. 123--293. isbn: 978-3-319-28736-2. doi: 10.1007/978-3-319-28736-2\_5.\par
\noindent\hangindent=2.2em\hangafter=1 [311]\ R. Scognamillo. “An elementary approach to the abelianization of the Hitchin system for arbitrary reductive groups.” Compos. Math. 110.1 (1998), pp. 17--37. issn: 0010-437X. doi: 10.1023/A:1000235107340.\par
\noindent\hangindent=2.2em\hangafter=1 [312]\ N. Seiberg. “Electric - magnetic duality in supersymmetric nonAbelian gauge theories.” Nucl. Phys. B 435 (1995), pp. 129--146. doi: 10.1016/0550-3213(94)00023-8.arXiv: hep-th/9411149.\par
\noindent\hangindent=2.2em\hangafter=1 [313]\ N. Seiberg and E. Witten. “Electric - magnetic duality, monopole condensation, and confinement in N=2 supersymmetric Yang-Mills theory.” Nucl. Phys. B 426 (1994). [Erratum: Nucl.Phys.B 430, 485--486 (1994)], pp. 19--52. doi: 10.1016/0550-3213(94)90124-4.arXiv: hep-th/9407087.\par
\noindent\hangindent=2.2em\hangafter=1 [314]\ N. Seiberg and E. Witten. “Monopoles, duality and chiral symmetry breaking in N=2 supersymmetric QCD.” Nucl. Phys. B 431 (1994), pp. 484--550. doi: 10.1016/0550-3213(94)90214-3.arXiv: hep-th/9408099.\par
\noindent\hangindent=2.2em\hangafter=1 [315]\ N. Seiberg and E. Witten. “Comments on string dynamics in six-dimensions.” Nucl. Phys. B 471 (1996), pp. 121--134. doi: 10.1016/0550-3213(96)00189-7.arXiv: hep-th/9603003.\par
\noindent\hangindent=2.2em\hangafter=1 [316]\ N. Seiberg. “Supersymmetry and Nonperturbative beta Functions.” Phys. Lett. B 206 (1988), pp. 75--80. doi: 10.1016/0370-2693(88)91265-8.\par
\noindent\hangindent=2.2em\hangafter=1 [317]\ N. Seiberg. “Naturalness versus supersymmetric nonrenormalization theorems.” Phys. Lett. B 318 (1993), pp. 469--475. doi: 10.1016/0370-2693(93)91541-T.arXiv: hep-ph/9309335.\par
\noindent\hangindent=2.2em\hangafter=1 [318]\ N. Seiberg. “Exact results on the space of vacua of four-dimensional SUSY gauge theories.” Phys. Rev. D 49 (1994), pp. 6857--6863. doi: 10.1103/PhysRevD.49.6857.arXiv: hep-th/9402044.\par
\noindent\hangindent=2.2em\hangafter=1 [319]\ N. Seiberg. “The Power of holomorphy: Exact results in 4-D SUSY field theories.” Particles, Strings, and Cosmology (PASCOS 94). May 1994, pp. 0357--369. arXiv: hepth/9408013.\par
\noindent\hangindent=2.2em\hangafter=1 [320]\ N. Seiberg. “New theories in six-dimensions and matrix description of M theory on $T^5$ and $T^5$ / Z(2).” Phys. Lett. B 408 (1997), pp. 98--104. doi: 10.1016/S0370-2693(97)00805-8.arXiv: hep-th/9705221.\par
\noindent\hangindent=2.2em\hangafter=1 [321]\ N. Seiberg. “Notes on theories with 16 supercharges.” Nucl. Phys. B Proc. Suppl. 67 (1998). Ed. by E. Gava, B. Greene, J. Louis, K. S. Narain, and S. Randjbar-Daemi, pp. 158--171. doi: 10.1016/S0920-5632(98)00128-5.arXiv: hep-th/9705117.\par
\noindent\hangindent=2.2em\hangafter=1 [322]\ N. Seiberg and E. Witten. “Gauge dynamics and compactification to three-dimensions.” Conference on the Mathematical Beauty of Physics (In Memory of C. Itzykson). June 1996, pp. 333--366. arXiv: hep-th/9607163.\par
\noindent\hangindent=2.2em\hangafter=1 [323]\ A. Sen. “Dyon - monopole bound states, selfdual harmonic forms on the multi - monopole moduli space, and SL(2,Z) invariance in string theory.” Phys. Lett. B 329 (1994), pp. 217-- 221. doi: 10.1016/0370-2693(94)90763-3.arXiv: hep-th/9402032.\par
\noindent\hangindent=2.2em\hangafter=1 [324]\ A. Sen. “Strong - weak coupling duality in four-dimensional string theory.” Int. J. Mod. Phys. A 9 (1994), pp. 3707--3750. doi: 10.1142/S0217751X94001497.arXiv: hepth/9402002.\par
\noindent\hangindent=2.2em\hangafter=1 [325]\ M. Serone. “Lectures on Resurgence in Integrable Field Theories” (May 2024). arXiv: 2405.02224 [hep-th].\par
\noindent\hangindent=2.2em\hangafter=1 [326]\ S. Shadchin. “Status report on the instanton counting.” SIGMA 2 (2006), p. 008. arXiv: hep-th/0601167.\par
\noindent\hangindent=2.2em\hangafter=1 [327]\ S.-H. Shao. “What’s Done Cannot Be Undone: TASI Lectures on Non-Invertible Symmetries” (Aug. 2023). arXiv: 2308.00747 [hep-th].\par
\noindent\hangindent=2.2em\hangafter=1 [328]\ A. D. Shapere and C. Vafa. “BPS structure of Argyres-Douglas superconformal theories” (Oct. 1999). arXiv: hep-th/9910182.\par
\noindent\hangindent=2.2em\hangafter=1 [329]\ M. A. Shifman. “Nonperturbative dynamics in supersymmetric gauge theories.” Prog. Part. Nucl. Phys. 39 (1997), pp. 1--116. doi: 10.1016/S0146-6410(97)00042-2.arXiv: hep-th/9704114.\par
\noindent\hangindent=2.2em\hangafter=1 [330]\ A. S. Sikora. “Character varieties.” Trans. Am. Math. Soc. 364.10 (2012), pp. 5173--5208. issn: 0002-9947. doi: 10.1090/S0002-9947-2012-05448-1.\par
\noindent\hangindent=2.2em\hangafter=1 [331]\ C. T. Simpson. “Higgs bundles and local systems.” Publ. Math., Inst. Hautes Étud. Sci. 75 (1992), pp. 5--95. issn: 0073-8301. doi: 10.1007/BF02699491.\par
\noindent\hangindent=2.2em\hangafter=1 [332]\ J. Song, Y. Cheng, B. Li, and B. Xu. “Drawing cone spherical metrics via Strebel differentials.” Int. Math. Res. Not. 2020.11 (2020), pp. 3341--3363. issn: 1073-7928. doi: 10.1093/imrn/rny103.\par
\noindent\hangindent=2.2em\hangafter=1 [333]\ A. M. Steane. “An introduction to spinors” (Dec. 2013). arXiv: 1312.3824 [math-ph].\par
\noindent\hangindent=2.2em\hangafter=1 [334]\ R. Steinberg. “The isomorphism and isogeny theorems for reductive algebraic groups.” J. Algebra 216.1 (1999), pp. 366--383. issn: 0021-8693. doi: 10.1006/jabr.1998.7776.\par
\noindent\hangindent=2.2em\hangafter=1 [335]\ M. J. Strassler. “An Unorthodox introduction to supersymmetric gauge theory.” Theoretical Advanced Study Institute in Elementary Particle Physics (TASI 2001): Strings, Branes and EXTRA Dimensions. Sept. 2003, pp. 561--638. doi: 10.1142/9789812702821\_0011.arXiv: hep-th/0309149.\par
\noindent\hangindent=2.2em\hangafter=1 [336]\ K. Strebel. Quadratic differentials. Vol. 5. Ergeb. Math. Grenzgeb., 3. Folge. Springer, Cham, 1984.\par
\noindent\hangindent=2.2em\hangafter=1 [337]\ A. Strominger and M. Dine. “Open p-branes.” Phys. Lett. B 383 (1996), pp. 44--47. doi: 10.1201/9781482268737-13.arXiv: hep-th/9512059.\par
\noindent\hangindent=2.2em\hangafter=1 [338]\ Y. Tachikawa. “N=2 S-duality via Outer-automorphism Twists.” J. Phys. A 44 (2011), p. 182001. doi: 10.1088/1751-8113/44/18/182001.arXiv: 1009.0339 [hep-th].\par
\noindent\hangindent=2.2em\hangafter=1 [339]\ Y. Tachikawa. N=2 supersymmetric dynamics for pedestrians. Dec. 2013. doi: 10.1007/978-3-319-08822-8.arXiv: 1312.2684 [hep-th].\par
\noindent\hangindent=2.2em\hangafter=1 [340]\ Y. Tachikawa. “A review on instanton counting and W-algebras.” New Dualities of Supersymmetric Gauge Theories. Ed. by J. Teschner. 2016, pp. 79--120. doi: 10.1007/978-3-319-18769-3\_4.arXiv: 1412.7121 [hep-th].\par
\noindent\hangindent=2.2em\hangafter=1 [341]\ J. Terning. Modern Supersymmetry: Dynamics and Duality. Oxford University Press, Dec. 2005. isbn: 9780198567639. doi: 10.1093/acprof:oso/9780198567639.001.0001.\par
\noindent\hangindent=2.2em\hangafter=1 [342]\ J. Teschner. “A guide to two-dimensional conformal field theory.” Integrability: from statistical systems to gauge theory. Lecture notes of the Les Houches summer school. Volume 106, Les Houches, France, June 6 -- July 1, 2016. Oxford: Oxford University Press, 2019, pp. 60--120. isbn: 978-0-19-882815-0. doi: 10.1093/oso/9780198828150.003.0002.\par
\noindent\hangindent=2.2em\hangafter=1 [343]\ A. Thomas. “A Gentle Introduction to the Non-Abelian Hodge Correspondence.” L’Enseignement Mathématique (Aug. 2024). published online first. arXiv: 2208.05940 [math.DG].\par
\noindent\hangindent=2.2em\hangafter=1 [344]\ W. P. Thurston. Collected works of William P. Thurston with commentary: IV. The geometry and topology of three-manifolds: with a preface by Steven P. Kerckhoff. Edited by Benson Farb, David Gabai and Steven P. Kerckhoff. Ed. by S. P. Kerckhoff, B. Farb, and D. Gabai. Providence, RI: American Mathematical Society (AMS), 2022. isbn: 978-1- 4704-5164-6.\par
\noindent\hangindent=2.2em\hangafter=1 [345]\ D. Tong. “TASI lectures on solitons: Instantons, monopoles, vortices and kinks.” Theoretical Advanced Study Institute in Elementary Particle Physics: Many Dimensions of String Theory. June 2005. arXiv: hep-th/0509216.\par
\noindent\hangindent=2.2em\hangafter=1 [346]\ D. Tong. Gauge Theory. DAMTP Cambridge Lecture Notes. 2018.\par
\noindent\hangindent=2.2em\hangafter=1 [347]\ P. K. Townsend and P. V. Landshoff. “D-branes from M-branes.” Phys. Lett. B 373 (1996), pp. 68--75. doi: 10.1201/9781482268737-14.arXiv: hep-th/9512062.\par
\noindent\hangindent=2.2em\hangafter=1 [348]\ C. Vafa and E. Witten. “A Strong coupling test of S duality.” Nucl. Phys. B 431 (1994), pp. 3--77. doi: 10.1016/0550-3213(94)90097-3.arXiv: hep-th/9408074.\par
\noindent\hangindent=2.2em\hangafter=1 [349]\ V. S. Varadarajan. Supersymmetry for mathematicians: an introduction. Vol. 11. Courant Lect. Notes Math. Providence, RI: American Mathematical Society (AMS); New York, NY: Courant Institute of Mathematical Sciences, 2004. isbn: 0-8218-3574-2.\par
\noindent\hangindent=2.2em\hangafter=1 [350]\ E. Vassiliou. “Flat bundles and holonomy homomorphisms.” Manuscr. Math. 42 (1983), pp. 161--170. issn: 0025-2611. doi: 10.1007/BF01169581.\par
\noindent\hangindent=2.2em\hangafter=1 [351]\ E. P. Verlinde. “Global aspects of electric - magnetic duality.” Nucl. Phys. B 455 (1995), pp. 211--228. doi: 10.1016/0550-3213(95)00431-Q.arXiv: hep-th/9506011.\par
\noindent\hangindent=2.2em\hangafter=1 [352]\ A. Voros. “The return of the quartic oscillator. The complex WKB method.” Ann. Inst. Henri Poincaré, Nouv. Sér., Sect. A 39 (1983), pp. 211--338. issn: 0020-2339.\par
\noindent\hangindent=2.2em\hangafter=1 [353]\ R. S. Ward. “On Selfdual gauge fields.” Phys. Lett. A 61 (1977), pp. 81--82. doi: 10.1016/0375-9601(77)90842-8.\par
\noindent\hangindent=2.2em\hangafter=1 [354]\ E. J. Weinberg and P. Yi. “Magnetic Monopole Dynamics, Supersymmetry, and Duality.” Phys. Rept. 438 (2007), pp. 65--236. doi: 10.1016/j.physrep.2006.11.002.arXiv: hep-th/0609055.\par
\noindent\hangindent=2.2em\hangafter=1 [355]\ S. Weinberg. The quantum theory of fields. Vol. 3: Supersymmetry. Cambridge University Press, June 2013. isbn: 978-0-521-67055-5.\par
\noindent\hangindent=2.2em\hangafter=1 [356]\ R. O. j. Wells. Differential analysis on complex manifolds. With a new appendix by Oscar Garcia-Prada. 3rd ed. Vol. 65. Grad. Texts Math. New York, NY: Springer, 2008. isbn: 978-0-387-73891-8. doi: 10.1007/978-0-387-73892-5.\par
\noindent\hangindent=2.2em\hangafter=1 [357]\ G. Wentzel. “Eine Verallgemeinerung der Quantenbedingungen für die Zwecke der Wellenmechanik.” Z. Phys. 38 (1926), pp. 518--529. issn: 0939-7922.\par
\noindent\hangindent=2.2em\hangafter=1 [358]\ J. Wess and J. Bagger. Supersymmetry and supergravity. Princeton, NJ, USA: Princeton University Press, 1992. isbn: 978-0-691-02530-8.\par
\noindent\hangindent=2.2em\hangafter=1 [359]\ J. Wess and B. Zumino. “Supergauge Transformations in Four-Dimensions.” Nucl. Phys. B 70 (1974). Ed. by A. Salam and E. Sezgin, pp. 39--50. doi: 10.1016/0550-3213(74)90355-1.\par
\noindent\hangindent=2.2em\hangafter=1 [360]\ A. M. Whitney. “A reduction theorem for totally positive matrices.” J. Anal. Math. 2 (1952), pp. 88--92. issn: 0021-7670. doi: 10.1007/BF02786969.\par
\noindent\hangindent=2.2em\hangafter=1 [361]\ E. Witten. “Dyons of Charge e theta/2 pi.” Phys. Lett. B 86 (1979), pp. 283--287. doi: 10.1016/0370-2693(79)90838-4.\par
\noindent\hangindent=2.2em\hangafter=1 [362]\ E. Witten. “Dynamical Breaking of Supersymmetry.” Nucl. Phys. B 188 (1981), p. 513. doi: 10.1016/0550-3213(81)90006-7.\par
\noindent\hangindent=2.2em\hangafter=1 [363]\ E. Witten. “An SU(2) Anomaly.” Phys. Lett. B 117 (1982). Ed. by M. A. Shifman, pp. 324-- 328. doi: 10.1016/0370-2693(82)90728-6.\par
\noindent\hangindent=2.2em\hangafter=1 [364]\ E. Witten. “Constraints on Supersymmetry Breaking.” Nucl. Phys. B 202 (1982), p. 253. doi: 10.1016/0550-3213(82)90071-2.\par
\noindent\hangindent=2.2em\hangafter=1 [365]\ E. Witten. “Some comments on string dynamics.” STRINGS 95: Future Perspectives in String Theory. July 1995, pp. 501--523. arXiv: hep-th/9507121.\par
\noindent\hangindent=2.2em\hangafter=1 [366]\ E. Witten. “Solutions of four-dimensional field theories via M-theory.” Nucl. Phys. B 500 (1997), pp. 3--42. doi: 10.1201/9781482268737-38.arXiv: hep-th/9703166.\par
\noindent\hangindent=2.2em\hangafter=1 [367]\ E. Witten. “Dynamics of quantum field theory.” Quantum fields and strings: a course for mathematicians. Vols. 1, 2. Material from the Special Year on Quantum Field Theory held at the Institute for Advanced Study, Princeton, NJ, 1996--1997. Providence, RI: AMS, American Mathematical Society, 1999, pp. 1119--1424. isbn: 0-8218-1987-9.\par
\noindent\hangindent=2.2em\hangafter=1 [368]\ E. Witten. “Supersymmetric index in four-dimensional gauge theories.” Adv. Theor. Math. Phys. 5 (2002), pp. 841--907. doi: 10.4310/ATMP.2001.v5.n5.a1.arXiv: hepth/0006010.\par
\noindent\hangindent=2.2em\hangafter=1 [369]\ E. Witten and D. I. Olive. “Supersymmetry Algebras That Include Topological Charges.” Phys. Lett. B 78 (1978), pp. 97--101. doi: 10.1016/0370-2693(78)90357-X.\par
\noindent\hangindent=2.2em\hangafter=1 [370]\ M. Wolf. “The Teichmüller theory of harmonic maps.” J. Differ. Geom. 29.2 (1989), pp. 449--479. issn: 0022-040X. doi: 10.4310/jdg/1214442885.\par
\noindent\hangindent=2.2em\hangafter=1 [371]\ D. Xie. “Higher laminations, webs and N=2 line operators” (Apr. 2013). arXiv: 1304.2390 [hep-th].\par
\noindent\hangindent=2.2em\hangafter=1 [372]\ C.-N. Yang and R. L. Mills. “Conservation of Isotopic Spin and Isotopic Gauge Invariance.” Phys. Rev. 96 (1954). Ed. by J.-P. Hsu and D. Fine, pp. 191--195. doi: 10.1103/PhysRev.96.191.\par
\noindent\hangindent=2.2em\hangafter=1 [373]\ A. B. Zamolodchikov. “Conformal Scalar Field on the Hyperelliptic Curve and Critical Ashkin-teller Multipoint Correlation Functions.” Nucl. Phys. B 285 (1987), pp. 481--503. doi: 10.1016/0550-3213(87)90350-6.\par
\noindent\hangindent=2.2em\hangafter=1 [374]\ A. B. Zamolodchikov. “Conformal symmetry in two-dimensional space: Recursion representation of conformal block.” Theor. Math. Phys. 73.1 (1987), pp. 1088--1093. doi: 10.1007/BF01022967.\par
\endgroup
